% =====================================================================
%  RAPPORT PFA - EMSI 2025-2026
%  Wydad Digital
%  Periode du projet : 1er aout 2026 - 31 aout 2026
%  Auteur : Ilyas Ait Maina
% =====================================================================

\documentclass[a4paper,12pt,oneside]{report}

% ----- Encodage & langue -----
\usepackage[utf8]{inputenc}
\usepackage[T1]{fontenc}
\usepackage[french]{babel}
\usepackage{lmodern}

% ----- Mise en page -----
\usepackage[top=2.8cm,bottom=2.5cm,left=2.5cm,right=2.5cm,
            headheight=15pt,headsep=20pt,footskip=28pt]{geometry}
\usepackage{setspace}
\onehalfspacing
\usepackage{indentfirst}
\usepackage{microtype}
\sloppy
\frenchspacing

% ----- Césure -----
\hyphenation{microservices Spring Boot 3.3, architecture, conteneurisation, redemarrage, disponibilite, maintenabilite, scalabilite, tracabilite, orchestree, persistante, predisposition, implementation, authentification, identification, recuperer, configuration, application, environnement, developpement, deploiement, integration, communication, messagerie, actualisation, conditionnement, conditionne, professionnalisation, inter-services, virtualisation, hierarchisation}
\lefthyphenmin=3
\righthyphenmin=4
\tolerance=2000
\emergencystretch=3em

% ----- Macros techniques -----
\newcommand{\microservice}[1]{\mbox{#1}}
\newcommand{\code}[1]{\mbox{\texttt{#1}}}
\newcommand{\API}[1]{\mbox{\textsc{#1}}}

% ----- Polices -----
\usepackage{mathptmx}
\usepackage[scaled=0.92]{helvet}
\renewcommand{\familydefault}{\rmdefault}

% ----- Mesures LaTeX (pour macro image) -----
\usepackage{calc}

% ----- Couleurs WAC -----
\usepackage{xcolor}
\definecolor{wacRouge}{HTML}{C8102E}
\definecolor{wacRougeFonce}{HTML}{8B0000}
\definecolor{wacOr}{HTML}{D4AF37}
\definecolor{wacGris}{HTML}{2C2C2C}
\definecolor{wacGrisClair}{HTML}{F5F5F5}
\definecolor{wacBlanc}{HTML}{FFFFFF}
\definecolor{wacBleu}{HTML}{1A3A6E}

% ----- Graphiques -----
\usepackage{graphicx}
% Les diagrammes PlantUML sont dans le MÊME dossier que ce .tex (Overleaf).
% \graphicspath permet d'écrire juste \insertimage{nom.png}{legende} sans préfixe.
\graphicspath{{./}{./images/}}
\usepackage{float}
\usepackage{caption}
\usepackage{tcolorbox}
\tcbuselibrary{skins, breakable}

% ----- Format paysage -----
\usepackage{pdflscape}

% ----- Tableaux -----
\usepackage{array}
\usepackage{booktabs}
\usepackage{longtable}
\usepackage{tabularx}
\usepackage{multirow}
\usepackage{makecell}
\usepackage{colortbl}

% ----- Listes -----
\usepackage{enumitem}
\setlist[itemize]{leftmargin=*,topsep=3pt,itemsep=1pt,parsep=1pt}
\setlist[enumerate]{leftmargin=*,topsep=3pt,itemsep=1pt,parsep=1pt}

% ----- En-tetes / pieds de page -----
\usepackage{fancyhdr}

% ----- Table des matieres -----
\usepackage{tocloft}
\renewcommand{\cftchapfont}{\color{wacRouge}\bfseries\large}
\renewcommand{\cftsecfont}{\color{wacGris}\bfseries}
\renewcommand{\cftsubsecfont}{\color{wacGris}}
\renewcommand{\cftdotsep}{1.5}
\setlength{\cftbeforechapskip}{6pt}

% ----- Liens hypertextes -----
\usepackage[colorlinks=true,
            linkcolor=wacRouge,
            urlcolor=wacBleu,
            citecolor=wacRougeFonce,
            pdfborder={0 0 0}]{hyperref}

% ----- Compteurs -----
\usepackage{chngcntr}
\counterwithout{figure}{chapter}
\counterwithout{table}{chapter}

% ----- Legendes -----
\captionsetup{font=small,labelfont={bf,color=wacRouge},labelsep=period,
              format=plain,justification=justified,
              singlelinecheck=false}
\renewcommand{\thefigure}{\arabic{figure}}
\renewcommand{\thetable}{\arabic{table}}
\renewcommand{\figurename}{Figure}
\renewcommand{\tablename}{Tableau}

% ----- En-tetes / pieds de page -----
\pagestyle{fancy}
\fancyhf{}
\fancyhead[L]{\color{wacRouge}\small\itshape Wydad Digital}
\fancyhead[R]{\color{wacGris}\small\itshape PFA 2025-2026}
\fancyfoot[C]{\color{wacGris}\small\thepage}
\renewcommand{\headrulewidth}{0.4pt}
\renewcommand{\headrule}{\hbox to\headwidth{\color{wacRouge}\leaders\hrule height \headrulewidth\hfill}}
\renewcommand{\footrulewidth}{0pt}

\fancypagestyle{plain}{%
  \fancyhf{}%
  \fancyfoot[C]{\color{wacGris}\small\thepage}%
  \renewcommand{headrulewidth}{0pt}%
}

% ----- Macros -----
\newcommand{\HRule}{\color{wacRouge}\rule{\linewidth}{1mm}}

\newtcolorbox{infobox}{
  colback=wacGrisClair,
  colframe=wacBleu,
  fonttitle=\bfseries\color{wacBlanc},
  title={A savoir},
  coltitle=wacBlanc,
  arc=2mm,
  boxrule=0.8pt,
  left=5pt, right=5pt, top=3pt, bottom=3pt,
  breakable
}

\newtcolorbox{keypointbox}{
  colback=wacGrisClair,
  colframe=wacRouge,
  fonttitle=\bfseries\color{wacBlanc},
  title={Point cle du projet},
  coltitle=wacBlanc,
  arc=2mm,
  boxrule=0.8pt,
  left=5pt, right=5pt, top=3pt, bottom=3pt,
  breakable
}

% Macro image : insere l'image si presente, sinon un cadre discret
% Usage : \insertimage{nom.png}{legende}        -> largeur 0.85\textwidth
%        \insertimage[0.95]{nom.png}{legende}  -> largeur personnalisee
%        \insertimage*[nom.png}{legende}        -> mode "fit" force (maxheight=0.7\textheight)
% MODE CASTRO : affiche le placeholder [nom.png] + une explication textuelle
%               au lieu de l'image raster. Le nom suffit, le detail est dans le texte.
\makeatletter
\newcommand{\insertimage}[2]{%
  \begin{figure}[H]\centering
    \fcolorbox{wacRouge}{wacGrisClair}{%
      \parbox[c][1.8cm][c]{0.70\textwidth}{\centering\itshape
        \color{wacRouge}\texttt{[#1]}}%
    }\par\vspace{0.2cm}
    {\small\itshape\color{wacGris}#2}%
    \label{fig:#1}%
  \end{figure}%
}
\makeatother

% Macro pratique : force l'image a tenir dans la page
% (largeur max 0.85\textwidth, hauteur max 0.70\textheight)
% MODE CASTRO : idem, placeholder + legende.
\makeatletter
\newcommand{\insertfit}[2]{%
  \begin{figure}[H]\centering
    \fcolorbox{wacRouge}{wacGrisClair}{%
      \parbox[c][1.8cm][c]{0.85\textwidth}{\centering\itshape
        \color{wacRouge}\texttt{[#1]}}%
    }\par\vspace{0.2cm}
    {\small\itshape\color{wacGris}#2}%
    \label{fig:#1}%
  \end{figure}%
}
\makeatother

% Macro compacte : image avec legende inline (pour les diagrammes)
% MODE CASTRO : idem, placeholder + legende.
\makeatletter
\newcommand{\insertimg}[2]{%
  \begin{figure}[H]\centering
    \fcolorbox{wacRouge}{wacGrisClair}{%
      \parbox[c][1.5cm][c]{0.75\textwidth}{\centering\itshape
        \color{wacRouge}\texttt{[#1]}}%
    }\par\vspace{0.15cm}
    {\small\itshape\color{wacGris}#2}%
  \end{figure}%
}
\makeatother

% Bandeau chapitre
\newcommand{\chapterbanner}[1]{%
  \clearpage
  \begin{center}
  \begin{tcolorbox}[
    colback=wacRouge,
    colframe=wacRouge,
    coltext=wacBlanc,
    arc=3mm,
    boxrule=0pt,
    width=0.95\textwidth,
    halign=center,
    left=4pt,right=4pt,top=5pt,bottom=5pt
  ]
  {\Large\bfseries\color{wacBlanc} #1}
  \end{tcolorbox}
  \end{center}
}

\begin{document}

\hyphenpenalty=10000
\exhyphenpenalty=10000

% =====================================================================
% PAGE DE GARDE
% =====================================================================
\thispagestyle{empty}
\begin{titlepage}
\begin{center}

\IfFileExists{logo_emsi.png}{%
  \includegraphics[width=4cm]{logo_emsi.png}\\[0.5cm]
}{%
  \fcolorbox{wacRouge}{wacGrisClair}{%
    \parbox[c][1.8cm][c]{4cm}{\centering\itshape\color{wacRouge}[Logo EMSI]%
  }}\\[0.5cm]
}

{\Huge\bfseries\color{wacRouge} EMSI}\\[0.2cm]
{\large Ecole Marocaine des Sciences de l'Ingenieur}\\[0.2cm]
{\small Cycle Ingenieur --- Annee universitaire 2025 / 2026}\\[1cm]

\color{wacRouge}\rule{\linewidth}{1.5mm}\\[0.5cm]

{\Huge\bfseries\color{wacGris} Wydad Digital}\\[0.3cm]
{\Large\itshape\color{wacRouge} Plateforme communautaire microservices\\
au service du Wydad Athletic Club}\\[0.5cm]

\color{wacRouge}\rule{\linewidth}{1.5mm}\\[1.2cm]

\begin{tcolorbox}[
  colback=wacGrisClair,
  colframe=wacRouge,
  arc=3mm,
  boxrule=1pt,
  width=0.88\textwidth,
  left=8pt,right=8pt,top=8pt,bottom=8pt
]
\centering
\begin{tabular}{@{}>{\raggedright\arraybackslash}p{5cm} >{\raggedleft\arraybackslash}p{7cm}@{}}
\textit{\large\color{wacGris}Realise par :}        & {\large\bfseries\color{wacRouge} Ilyas Ait Maina} \\[0.3cm]
\textit{\large\color{wacGris}Encadrant pedagogique :} & {\large\bfseries\color{wacRouge} Dr. Saadia Drissi} \\[0.3cm]
\textit{\large\color{wacGris}Periode du projet :}     & {\large\bfseries\color{wacRouge} 1\textsuperscript{er} -- 31 aout 2026} \\
\end{tabular}
\end{tcolorbox}\\[0.8cm]

\begin{tcolorbox}[
  colback=wacBlanc,
  colframe=wacOr,
  arc=3mm,
  boxrule=0.8pt,
  width=0.88\textwidth,
  left=8pt,right=8pt,top=8pt,bottom=8pt
]
\centering
\textit{\small\color{wacGris}Membres du jury :}\\[0.4cm]
\color{wacGris}\rule{0.65\textwidth}{0.3pt} \\[0.25cm]
\rule{0.65\textwidth}{0.3pt} \\[0.25cm]
\rule{0.65\textwidth}{0.3pt} \\
\end{tcolorbox}

\vfill

{\small\color{wacGris} Soutenu le \rule{3cm}{0.3pt} a Casablanca devant le jury ci-dessus.}\\[1cm]

{\small\bfseries\color{wacRouge} EMSI} {\small\color{wacGris}--- Promotion 2025--2026}

\end{center}
\end{titlepage}

% =====================================================================
% PAGES LIMINAIRES
% =====================================================================
\pagenumbering{roman}
\setcounter{page}{1}

\clearpage
\chapter*{Dedicaces}
\addcontentsline{toc}{chapter}{Dedicaces}

\begin{center}
\itshape\large
\color{wacRouge}\rule{4cm}{0.5pt}\\[0.8cm]

{\Large\color{wacGris} A mes parents,}\\[0.3cm]
qui m'ont soutenu tout au long de ce parcours et sans qui rien de tout
cela n'aurait ete possible.\\[0.8cm]

{\Large\color{wacGris} A mes freres et soeurs,}\\[0.3cm]
pour leur patience et leur soutien de tous les instants.\\[0.8cm]

{\Large\color{wacGris} A mes amis et camarades de promotion,}\\[0.3cm]
avec qui j'ai partage les meilleurs moments a l'EMSI.\\[0.8cm]

{\Large\color{wacGris} A toute la communaute Wydad,}\\[0.3cm]
qui m'a inspire dans la conception de cette plateforme.\\[1.2cm]

{\Large\color{wacRouge} Je vous dedie ce rapport.}\\[0.5cm]

\color{wacRouge}\rule{4cm}{0.5pt}
\end{center}

\clearpage
\chapter*{Remerciements}
\addcontentsline{toc}{chapter}{Remerciements}

Avant de presenter ce travail, je tiens a remercier toutes les
personnes qui ont contribue, de pres ou de loin, a la realisation de
ce projet de fin d'annee.

Je remercie en premier lieu \textbf{\color{wacRouge}mon encadrant
pedagogique} pour sa disponibilite, ses orientations avisées et la
rigueur qu'il m'a transmise tout au long de ce travail. Ses conseils
methodologiques m'ont permis de structurer ma demarche et d'atteindre
les objectifs fixes dans le calendrier serre du mois d'aout 2026.

Je remercie egalement l'equipe pedagogique de l'\textbf{EMSI} pour la
qualite de la formation dispensée. Les connaissances acquises en
architecture logicielle, en genie logiciel, en bases de donnees et
en genie DevOps ont constitue le socle technique sur lequel s'est
construit ce projet.

Un grand merci a mes collegues de promotion avec qui j'ai pu
echanger, apprendre et progresser. Les discussions techniques, les
relectures croisees et l'entraide quotidienne ont considerablement
contribue a l'aboutissement du projet.

Enfin, je remercie ma famille pour son soutien moral indefectible
durant toute la periode de conception, de developpement et de
redaction de ce rapport.

\clearpage
\renewcommand{\contentsname}{\color{wacRouge}Table des matieres}
\tableofcontents

\clearpage
\renewcommand{\listfigurename}{\color{wacRouge}Liste des figures}
\listoffigures
\addcontentsline{toc}{chapter}{Liste des figures}

\clearpage
\renewcommand{\listtablename}{\color{wacRouge}Liste des tableaux}
\listoftables
\addcontentsline{toc}{chapter}{Liste des tableaux}

% =====================================================================
% CORPS
% =====================================================================
\pagenumbering{arabic}
\setcounter{page}{1}

% =====================================================================
% INTRODUCTION GENERALE
% =====================================================================
\chapterbanner{INTRODUCTION GENERALE}
\chapter*{Introduction generale}
\addcontentsline{toc}{chapter}{Introduction generale}

\begin{keypointbox}
Le projet \textbf{Wydad Digital} porte sur la conception, le
developpement et le deploiement d'une plateforme numerique
microservices pour le Wydad Athletic Club. Il s'articule autour d'un
monorepo Maven hebergeant 11 microservices Spring Boot 3.3, un frontal
Angular 19, une passerelle Spring Cloud Gateway, plusieurs bases
PostgreSQL separees, et plusieurs services tiers (Cloudinary, LiveKit).
Le tout est conteneurise via Docker Compose, deploye sur une VM
Azure unique, et orchestree par un pipeline CI/CD GitHub Actions.
La periode de realisation s'etend du \textbf{1\textsuperscript{er}
aout 2026 au 31 aout 2026}.
\end{keypointbox}

\section{Contexte general}

Le Wydad Athletic Club (WAC), fonde en \textbf{\color{wacRouge}1937} a
Casablanca, est l'un des clubs sportifs les plus emblematiques du
Maroc et du continent africain. Avec un palmares exceptionnel
(Botola Pro, Ligue des Champions de la CAF remportee en 1992, 2017 et
2022), le club dispose d'une communaute de supporters tres
importante, repartie au Maroc, en Afrique et dans la diaspora.

Cette communaute interagit aujourd'hui avec le club via un patchwork
de canaux non coordonnes : site web institutionnel statique, pages
sur les reseaux sociaux, contacts telephoniques, convocations diffusees
par WhatsApp, et billetterie physique aux guichets du stade. Cette
fragmentation genere de nombreux points de friction : absence
d'adhesion en ligne, pas de billetterie digitale, gestion
administrative manuelle, absence de canal de communication~interne
tracable, impossibilite d'organiser des elections securisees en
ligne.

L'idee de \textbf{Wydad Digital}, nee dans le cadre academique
encadre de l'EMSI en partenariat avec le WAC, est de doter le club
d'une plateforme numerique centralisee capable d'adresser
simultanement l'ensemble de ces besoins.

\section{Problematique}

\begin{center}
\begin{tcolorbox}[
  colback=wacGrisClair,
  colframe=wacRouge,
  arc=3mm,
  boxrule=1pt,
  width=0.95\textwidth,
  left=8pt,right=8pt,top=6pt,bottom=6pt
]
\itshape\large\color{wacGris}
Comment concevoir et deployer, dans le respect des contraintes
budgetaires d'un club sportif, une plateforme numerique complete,
securisee et modulaire qui reponde simultanement aux besoins de
plusieurs categories d'utilisateurs --- visiteurs, adherents,
parents, joueurs, entraineurs, staff, journalistes et president de
club --- tout en couvrant l'integralite des processus metiers du
club (sportif,~commercial,~communication~interne,~gouvernance
participative, engagement communautaire) ?
\end{tcolorbox}
\end{center}

Cette problematique se decline en plusieurs questions secondaires :

\begin{itemize}
  \item Comment garantir la \textbf{\color{wacRouge}securite} des
        donnees personnelles et financieres (authentification forte,
        gestion fine des roles, communication inter-services protegee
        par secret partage) ?
  \item Comment assurer la \textbf{\color{wacRouge}scalabilite} et la
        \textbf{\color{wacRouge}maintenabilite} a long terme
        (architecture decouplee, code testable, documentation a
        jour) ?
  \item Comment proposer une \textbf{\color{wacRouge}experience
        utilisateur moderne} (front-end reactif, notifications temps
        reel, messagerie integree, appels video) ?
  \item Comment maitriser les \textbf{\color{wacRouge}couts
        d'hebergement} (deploiement sur une seule VM Azure, services
        tiers gratuits ou self-hosted) ?
\end{itemize}

\section{Objectifs du projet}

Le but du PFA est la \textbf{conception, le developpement et le
deploiement} d'une plateforme web microservices complete baptisee
\textbf{Wydad Digital}. Cette plateforme se doit d'etre :

\begin{enumerate}
  \item \textbf{\color{wacRouge}Complete} : couvrir tous les metiers
        du club (sport, billetterie, abonnement, academie,
        gouvernance, communication, engagement).
  \item \textbf{\color{wacRouge}Securisee par defaut} : authentification
        JWT, gestion fine des roles, validation serveur systematique,
        secret partage inter-services, propagation d'identite via
        en-tetes \code{X-User-*}, anti-IDOR strict.
  \item \textbf{\color{wacRouge}Modulaire} : architecture~microservices,
        un service par thematique~metier, communication~inter-services
        via REST protegee par \code{X-Internal-Secret}.
  \item \textbf{\color{wacRouge}Pereenne et economique} : deploiement
        sur une VM unique Azure Standard\_B2s\_v2 (2 vCPU, 8~Go RAM,
        30~Go SSD) en region Spain Central.
  \item \textbf{\color{wacRouge}Internationalisable et~responsive} :
        interface en francais, prete pour l'arabe et l'anglais,
        adaptation mobile/tablette/desktop.
\end{enumerate}

\section{Methode de travail}

Pour mener a bien ce projet dans la periode 1\textsuperscript{er}
-- 31 aout 2026 (soit un mois en parallele de l'enseignement), nous
avons opte pour une demarche structuree en plusieurs vagues
complementaires, dans l'esprit \textbf{Agile Scrum adapte a un projet
en solo} :

\begin{enumerate}
  \item \textbf{Etude de l'existant et recueil des besoins} : analyse
        des processus manuels, identification des points de friction,
        entretiens avec le porteur de projet.
  \item \textbf{Specification fonctionnelle et non fonctionnelle} :
        acteurs, cas d'utilisation, exigences de securite, performance,
        disponibilite, maintenabilite.
  \item \textbf{Conception} : choix de l'architecture~microservices,
        modelisation des flux (cas d'utilisation, classes, sequence,
        etat, activite, collaboration), modele de donnees, dictionnaire
        des donnees.
  \item \textbf{Developpement iteratif en vagues} :
        \begin{itemize}
          \item \textbf{Vague 0} : socle des roles et statuts de
                comptes.
          \item \textbf{Vague 1} : authentification et abonnements.
          \item \textbf{Vague 2} : billetterie VIP avec QR codes
                signes HMAC-SHA256.
          \item \textbf{Vague 3} : convocations, seances, mediateque.
          \item \textbf{Vague 4} : messagerie temps reel (WebSocket
                STOMP).
          \item \textbf{Vague 5} : appels video LiveKit.
          \item \textbf{Vague 6} : elections et sondages.
          \item \textbf{Vague 7 (B.18)} : parent supporter, achat
                pour enfant.
          \item \textbf{Vague 8 (C.21)} : espace president de club.
        \end{itemize}
  \item \textbf{Deploiement} sur VM Azure via Docker Compose, CI/CD
        GitHub Actions.
\end{enumerate}

\section{Plan du rapport}

Le rapport s'articule autour de quatre chapitres principaux :

\begin{itemize}
  \item \textbf{Chapitre 1} : Presentation du cadre du projet.
  \item \textbf{Chapitre 2} : Specification des besoins.
  \item \textbf{Chapitre 3} : Conception du systeme.
  \item \textbf{Chapitre 4} : Realisation du systeme.
\end{itemize}

% =====================================================================
% CHAPITRE 1
% =====================================================================
\chapterbanner{CHAPITRE 1 : PRESENTATION DU CADRE DU PROJET}
\chapter{Présentation du cadre du projet}

\section{Introduction}

Ce premier chapitre pose le contexte du projet. On y presente le
Wydad Athletic Club, l'etude de l'existant, la critique de cet
existant, la solution proposee, la methode de travail adoptee, ainsi
que le planning previsionnel de realisation.

\section{Présentation du Wydad Athletic Club}

\subsection{Historique}

Le Wydad Athletic Club a ete fonde le \textbf{\color{wacRouge}8 mai
1937} a Casablanca par cinq personnalites issues de la bourgeoisie
casablancaise. A l'origine simple club omnisports, le WAC s'est
progressivement impose comme l'une des institutions sportives
majeures du Royaume du Maroc.

\subsection{Palmares et rayonnement}

Le palmares du WAC est exceptionnel. En football, le club a remporte
a plus de 20 reprises le Championnat du Maroc (Botola Pro), une
dizaine de fois la Coupe du Trone, et surtout la \textbf{Ligue des
Champions de la CAF} a trois reprises (1992, 2017, 2022), faisant
du WAC l'un des clubs les plus titres du continent africain.

\subsection{Organisation interne}

Le WAC est structure en plusieurs sections sportives (football,
basketball, handball, athletisme, natation), une academie de
formation des jeunes, et plusieurs organes de gouvernance :
Bureau directeur, President, Assemblee generale, Commission
electorale, Staff sportif et Administration.

\section{Etude de l'existant}

\subsection{Site web officiel}

Le site officiel du club (\url{https://www.wac.ma/}) est un site
vitrine statique presentant l'historique, le palmares, l'effectif
professionnel, le calendrier des matchs et les actualites. Il n'est
pas interactif : aucun achat en ligne n'est possible.

\subsection{Reseaux sociaux}

Le club est present sur Facebook, X, Instagram, YouTube. Ces canaux
sont utilises principalement pour la communication~descendante et
la relation supporters, mais ils ne permettent aucune interaction
structuree (paiements, adhesions, votes, sondages).

\subsection{Gestion des adherents et billetterie}

La gestion des adherents est realisee a l'aide d'un tableur partage
et de contacts telephoniques. La billetterie physique est geree aux
guichets du stade le jour du match. Aucune billetterie en ligne
n'est proposee.

\subsection{Communication interne}

Les convocations, comptes-rendus de~match, et communications
internes circulent par e-mail, WhatsApp, ou de vive voix. Aucune
tracabilite ni archivage structure n'est garanti.

\section{Critique de l'existant}

L'analyse de l'existant met en evidence les limites suivantes :
fragmentation, manque d'interactivite, insecurite, absence de
tracabilite, scalabilite limitee.

\section{Solution proposee}

\subsection{Vision}

Pour repondre a ces limites, nous proposons la conception et le
deploiement d'une plateforme digitale complete, securisee, modulaire
et evolutive baptisee \textbf{\color{wacRouge}Wydad Digital}.

\subsection{Six piliers fonctionnels}

\begin{enumerate}
  \item \textbf{Adhesion et billetterie} : inscription multi-roles,
        abonnement saisonnier, billetterie par evenement, QR code,
        paiement en ligne (carte et E-Cash), cas B.18.
  \item \textbf{Academie et roles sportifs} : inscription des enfants,
        fiches roster, convocations, seances, fiches medicales,
        messagerie interne.
  \item \textbf{Gouvernance participative} : elections securisees,
        sondages, publication publique.
  \item \textbf{Communication officielle et~interne} :~actualites,
        medias, annonces, messagerie, chat temps reel, appels video.
  \item \textbf{Engagement communautaire} : points, badges,
        leaderboard.
  \item \textbf{Administration} : back-office complet.
\end{enumerate}

\subsection{Contraintes non fonctionnelles}

La solution doit respecter : securite (JWT, BCrypt, secret partage,
anti-IDOR), performance (P95 < 500~ms, 500~utilisateurs
concurrents), disponibilite (sauvegardes, redemarrage automatique),
maintenabilite (code documente, CI/CD), economie (cout hebergement
< 30~\$/mois).

\section{Methode de travail}

\subsection{Methode Agile Scrum}

Le projet est decoupe en vagues de developpement, chacune produisant
un increment fonctionnel teste et deployable. Product Owner, Scrum
Master et Equipe de developpement sont assures par le porteur de
projet et moi-meme.

\section{Planning previsionnel}

\begin{table}[H]
\centering
\small
\caption{Planning previsionnel du projet (aout 2026)}
\label{tab:planning}
\begin{tabular}{@{}p{5.2cm}p{5.5cm}p{1.8cm}p{2.2cm}@{}}
\toprule
\textbf{Phase} & \textbf{Description} & \textbf{Duree} & \textbf{Periode} \\
\midrule
Cadrage \& conception & Etude de l'existant, choix tech. & 1 sem. & 1-7 aout \\
Vague 0 & Socle roles/statuts, validation admin & 1 j. & 8 aout \\
Vague 1 & Authentification, JWT, KYC, abonnements & 2 j. & 9-10 aout \\
Vague 2 & Billetterie VIP, QR signes, anti-rejeu & 3 j. & 11-13 aout \\
Vague 3 & Convocations, seances, mediateque & 2 j. & 14-15 aout \\
Vague 4 & Messagerie WebSocket, STOMP & 3 j. & 16-18 aout \\
Vague 5 & Appels video LiveKit & 2 j. & 19-20 aout \\
Vague 6 & Elections et sondages & 2 j. & 21-22 aout \\
Vague 7 (B.18) & Parent supporter, achat pour fils & 2 j. & 23-24 aout \\
Vague 8 (C.21) & President de club, badge, annuaire & 2 j. & 25-26 aout \\
Deploiement final & VM Azure, CI/CD, soutenance & 5 j. & 27-31 aout \\
\bottomrule
\end{tabular}
\end{table}

\section{Conclusion}

Ce premier chapitre a permis de poser le contexte du projet et de
presenter la solution proposee. Le chapitre suivant sera consacre a
la specification detaillee des besoins.

% =====================================================================
% CHAPITRE 2
% =====================================================================
\chapterbanner{CHAPITRE 2 : SPECIFICATION DES BESOINS}
\chapter{Spécification des besoins}

\section{Introduction}

Ce chapitre formalise les besoins du systeme. On y decrit les acteurs,
les diagrammes de cas d'utilisation, les besoins fonctionnels et non
fonctionnels, et les principaux cas d'utilisation.

\section{Acteurs du systeme}

\begin{infobox}
Le systeme interagit avec \textbf{8 categories d'acteurs}, du simple
visiteur (sans compte) au president de club. Chaque acteur dispose
d'un role Spring Security specifique, valide cote serveur a chaque
requete via les en-tetes \code{X-User-*} propages par la passerelle.
\end{infobox}

\begin{table}[H]
\centering
\small
\caption{Acteurs du systeme et leurs roles}
\label{tab:acteurs}
\begin{tabular}{@{}p{4cm}p{2.6cm}p{8.4cm}@{}}
\toprule
\textbf{Acteur} & \textbf{Type} & \textbf{Role principal} \\
\midrule
Visiteur (non authentifie) & Humain externe & Consulter les pages publiques, acheter un billet en mode guest \\
Adherent / Supporter & Humain externe & Acheter billets, s'abonner, voter aux sondages, commenter \\
Parent & Humain externe & Inscrire un enfant, suivre ses entrainements, acheter pour ses enfants \\
Joueur & Humain interne & Consulter ses convocations, echanger avec l'encadrement \\
Entraineur / Staff & Humain interne & Gerer son equipe, planifier les seances, communiquer avec joueurs et parents \\
Journaliste & Humain externe & Demander une accreditation, couvrir les matchs, recevoir un badge officiel \\
President de club & Humain interne & Piloter la discipline, dialoguer avec joueurs et entraineurs, gerer recus et annonces \\
Administrateur & Humain interne & Valider les comptes, gerer les plans, superviser la plateforme \\
\bottomrule
\end{tabular}
\end{table}

\section{Besoins fonctionnels par pilier}

\subsection{Adhesion et billetterie}

L'inscription est multi-roles (VISITEUR, ADHERENT, PARENT, JOUEUR,
ENTRAINEUR, STAFF, JOURNALISTE, PRESIDENT) avec validation
administrative pour les roles privilegies. La connexion se fait par
e-mail et mot de passe, avec gestion de session JWT (Access Token
15 min, Refresh Token 7 jours). Le profil utilisateur gere la
photo, les informations personnelles, les preferences linguistiques
et l'avatar Cloudinary. La carte d'adherent numerique est disponible
en PDF + QR code telechargeable. L'abonnement saisonnier (carte
supporter annuel) beneficie d'un tarif adherent (reduction fidelite)
et les plans sont geres par l'administrateur.

Pour la billetterie, l'utilisateur selectionne la zone
(TRIBUNE\_OFFICIELLE, TRIBUNE\_HONNEUR, VIRAGE\_NORD, etc.) et le
nombre de places, procede au paiement (carte bancaire mockee ou
portefeuille E-Cash), et recoit un QR code unique signe HMAC-SHA256
verifie a l'entree du stade (anti-rejeu). Le billet PDF est
telechargeable avec sa facture associee.

Le cas B.18 merite qu'on s'y arrete : un parent doit pouvoir
acheter un billet ou un abonnement \emph{au nom de son enfant
academie}. Le systeme cree automatiquement un \emph{User shadow}
pour l'enfant (\code{enfant-\{academyId\}@wac.parent}) et le ticket
ou abonnement est rattache a cet enfant, avec tracabilite via les
colonnes \code{beneficiary\_academy\_member\_id} et
\code{parent\_payer\_email}. Le paiement reste debite sur le wallet
E-Cash du parent.

\subsection{Academie et roles sportifs}

L'inscription d'un enfant a l'academie suit une formule multi-etapes
(identite, sport, categorie, pieces justificatives). Les fiches
roster (joueurs + staff) sont gerees par discipline (FOOTBALL,
BASKETBALL, HANDBALL) et categorie (U15, U17, U18, U20, SENIOR).
L'entraineur planifie les seances d'entrainement et les convocations.
L'espace parent permet de suivre les enfants, de gerer les documents
justificatifs et de communiquer avec l'encadrement. L'espace joueur
permet de consulter les convocations, les seances, et d'echanger
avec le coach et les coequipiers. Une mediateque tactique permet
l'upload de videos, schemas et notes. Un emargement electronique
gere la liste des presents, absents et excuses.

\subsection{Gouvernance participative}

L'administrateur cree les scrutins (elections ou sondages) et definit
les candidats, dates d'ouverture et de fermeture. Le \textbf{vote a
l'election presidentielle est conditionne} (B.18) : seuls les membres
ayant un abonnement saisonnier ACTIF peuvent voter. Le backend
verifie via \code{AuthSubscriptionClient.isActiveSubscriber(email)},
et renvoie 403 + code \code{VOTE\_REQUIRES\_MEMBERSHIP} en cas
d'echec. ADMIN est exempte (override organisationnel). Les sondages
a choix multiples sont ouverts a tous les connectes, sans condition
d'abonnement. Le depouillement est automatique et la publication des
resultats est publique (sans authentification).

\subsection{Communication officielle et~interne}

La messagerie privee 1-a-1 (type WhatsApp) gere le statut en ligne
et la derniere connexion. Les chats de groupe (staff, equipe,
supporters, president de discipline) sont supportes. Les nouveaux
messages sont diffuses en temps reel via WebSocket STOMP + SockJS.
Les pieces jointes (images, PDF) sont uploadees sur Cloudinary avec
URL signees (TTL 1 heure). Le cas C.21 permet a un joueur d'ecrire
au president de sa discipline (et inversement) : le backend verifie
le role du destinataire via \code{AuthUserInfoClient} et autorise
l'envoi si les disciplines matchent. Les appels video et audio
passent par LiveKit (jeton JWT signe 2h, annulation, agenda
partage).

\subsection{Engagement communautaire}

Le systeme de points de fidelite recompense chaque action (achat,
vote, participation). Le leaderboard classe les supporters les plus
actifs. Des badges et niveaux de membre sont attribues
automatiquement.

\subsection{Administration}

Le back-office permet la validation des comptes en attente, la
gestion des plans d'abonnement, la gestion de la grille tarifaire
des sections de billetterie, la moderation des annonces et
sondages.

\subsection{Espace president de club (C.21)}

L'inscription PRESIDENT se fait avec choix de la discipline
(FOOTBALL, BASKETBALL, HANDBALL) et depot d'un justificatif
d'identite. La validation admin active le compte (\code{EN\_ATTENTE}
$\to$ \code{VALIDE}). Le dashboard president affiche un
\textbf{badge dore GOL} (au lieu du rouge club standard) avec nom,
discipline et statut. L'annuaire est filtre par discipline : le
president ne voit et ne peut parler qu'aux joueurs + entraineurs de
\emph{sa} discipline.

\section{Besoins non fonctionnels}

\subsection{Securite}

L'authentification est JWT (RFC 7519) avec algorithme HS256. Les
refresh tokens sont a rotation systematique. Les mots de passe sont
haches par BCrypt (cout 12). La validation cote serveur est
systematique (Jakarta Validation). Les QR codes de billetterie sont
signes HMAC-SHA256 (anti-falsification, anti-rejeu). Le secret
partage inter-services \code{X-Internal-Secret} est compare en
temps constant. L'identite est propagee par la passerelle via trois
en-tetes \code{X-User-Id}, \code{X-User-Email},
\code{X-User-Role}, lus par le \code{UserContextFilter} cote service.
L'anti-IDOR est strict.

\subsection{Performance}

Le temps de reponse P95 est inferieur a 500~ms pour 95~{\%} des
requetes. Le systeme supporte 500 utilisateurs concurrents.

\subsection{Disponibilite}

Le taux de disponibilite mensuel est superieur a 99,5~{\%}. Les
conteneurs Docker redemarrent automatiquement (restart policy). Des
sauvegardes quotidiennes \code{pg\_dumpall} sont effectuees.

\subsection{Maintenabilite}

Le code est documente, l'architecture en couches (Controller /
Service / Repository) est respectee, le pipeline CI/CD GitHub
Actions gere le build et le redeloiement, et les logs sont
structures.

\subsection{Compatibilite et portabilite}

L'application est accessible depuis les navigateurs modernes
(Chrome, Firefox, Edge, Safari). Elle est responsive (mobile,
tablette, desktop).

\section{Diagrammes de cas d'utilisation}

\subsection{Vues globales (par cote metier)}

\insertimage{uc-01a-vue-globale-supporter.png}{Diagramme de cas
d'utilisation (UC-01a, cote SUPPORTER) : Espace public et
externe --- VISITEUR, ADHERENT, PARENT, JOURNALISTE. Cette vue
presente les cas d'utilisation accessibles aux utilisateurs
externes (consulter le catalogue, acheter un billet, s'abonner,
demander une accreditation presse, suivre un enfant a l'academie,
participer a la gouvernance, etc.).}

\insertimage{uc-01b-vue-globale-club.png}{Diagramme de cas
d'utilisation (UC-01b, cote CLUB) : Entites internes du club ---
JOUEUR, ENTRAINEUR, STAFF, PRESIDENT, ADMIN. Cette vue presente
les cas d'utilisation des acteurs internes (convocations, seances,
feuille de match, recus, appels video, messagerie unifiee, etc.).}

\subsection{Detail par role interne (VAGUE 2)}

\begin{keypointbox}
Les diagrammes suivants ont ete \textbf{enrichis} a partir du code
source qui compile (controllers, services, DTOs et routes
Spring Boot 3.3 cote back ; composants standalone et services
Angular 19 cote front). Chaque diagramme suit la convention UML
2.5 : swimlanes par acteur, packages par sous-domaine fonctionnel,
relation \code{<<include>>} losange plein pour les dependances
obligatoires, heritage entre acteurs (ENTRAINEUR herite de STAFF).
Le respect strict de la separation des responsabilites (un UC par
action metier, pas de chevauchement entre packages) garantit la
coherence de la modelisation.
\end{keypointbox}

\insertimage{uc-02-supporter-public.png}{UC-02 --- SUPPORTER
public : VISITEUR, ADHERENT, PARENT, JOURNALISTE sans authentification
forte. Consulter actualites, palmares, classement, calendrier,
sondages publics, billetterie, plans abonnement, academies,
partenaires, stade, FAQ.}

\insertimage{uc-03-supporter-achat.png}{UC-03 --- SUPPORTER
achat : ADHERENT, PARENT achete billetterie, abonnement saisonnier,
boutique, E-Cash, dons, code promo. 5 cas d'utilisation repartis
en sous-packages.}

\insertimage{uc-04-supporter-academie.png}{UC-04 --- SUPPORTER
academie : PARENT gere ses enfants academie, paie pour eux
(B.18 User shadow + SportsAcademyClient + anti-IDOR strict).}

\insertimage{uc-05-supporter-gamification.png}{UC-05 --- SUPPORTER
gamification : ADHERENT consulte son profil gami, leaderboard,
pronostics (C1), badges, defis temporels (C2), points (C3).}

\insertimage{uc-06-supporter-votes.png}{UC-06 --- SUPPORTER
gouvernance : ADHERENT, PARENT votent aux elections (B.8.h
titulaire ACTIF + ADMIN exempt) et aux sondages. Resultats
publics sans auth apres publication.}

\insertimage{uc-07-club-joueur.png}{UC-07 --- JOUEUR : 25 cas
d'utilisation repartis en 8 packages. Auth, supporter (B.12
abonnement, B.28 billetterie), convocations (B.18.b nominatives),
feuille de match (B.18.b cochable par joueur), presence et
documents (KYC Phase 0), profil et stats, messagerie unifiee
(V2.5 team-chat photos + 1-to-1 persistant + canal president
+ selecteur libre 1..N).}

\insertimage{uc-08-club-staff-entraineur.png}{UC-08 --- STAFF et
ENTRAINEUR (heritage) : 24 cas d'utilisation, 8 packages.
Auth, annuaire discipline, convocations (creation + suivi
reponses), feuille de match (validation admin), seances,
documents KYC, messagerie unifiee V2.5 (memes sous-onglets que
joueur), comptes-rendus sportifs (C.21 recus de salaire).}

\insertimage{uc-09-club-president.png}{UC-09 --- PRESIDENT
(de SA discipline) : 22 cas d'utilisation, 8 packages.
Auth, annuaire de SA discipline, recus salaire/prime (C.21
avec IDOR proprio), documents president (brouillon + soumission
admin), messagerie unifiee V2.5, supporter (B.12 le president
peut acheter son propre abonnement).}

\insertimage{uc-10-club-messagerie.png}{UC-10 --- Messagerie
unifiee V2.5 (4 acteurs) : 16 cas d'utilisation transversaux.
Team-chat groupe (STOMP /topic + REST etendu pour photos),
messagerie 1-to-1 (avec photo), canal president (V2.5 nouveau),
selection libre 1..N destinataires (V2.5 nouveau), annonces
equipe, medias. Anti-IDOR strict sameGroup sportType+category
equalsIgnoreCase.}

\insertimage{uc-11-club-admin.png}{UC-11 --- ADMIN (back-office) :
42 cas d'utilisation, 7 packages. CRUD complet sur 9 roles
utilisateurs, validation comptes EN\_ATTENTE + KYC Phase 0,
CRUD contenus (actualites, mediatheque, matchs, classements,
seances, convocations, feuilles de match, effectif, fiches
publiques, palmares, legendes, sponsors, carrousels backgrounds,
reseaux sociaux, stade), billetterie + boutique + abonnements,
gamification (C3 bareme + C2 defis + badges), gouvernance
(elections, sondages), academy (B.18), presse (B.17 validation
demandes), recus salaire (C.21).}

\insertimage{uc-12-club-journaliste.png}{UC-12 --- JOURNALISTE
(B.17) : demande d'accreditation presse multi-matchs avec
multipart (photo + media + n° carte). 5 etats EN\_ATTENTE /
VALIDE / REFUSE. Badge PDF 19 Ko auto-genere (OpenPDF) sur
demande VALIDE. IDOR badge strict : self OU ADMIN.

\section{Cas d'utilisation principaux}

\subsection{CU-01 : Acheter un billet en ligne}

L'acteur principal est le Visiteur (sans compte) ou l'Adherent. La
pre-condition est qu'au moins un evenement ouvert a la billetterie
existe. Le scenario nominal est : (1) selection d'un evenement et
d'une section, (2) saisie des informations personnelles, (3) paiement,
(4) generation des QR codes signes et envoi du PDF. La post-condition
est la creation des billets, la signature du QR et l'envoi de
l'e-mail.

\subsection{CU-02 : S'authentifier et gerer sa session}

L'acteur principal est tout utilisateur enregistre. La pre-condition
est un compte existant avec statut \code{VALIDE}. Le scenario
nominal est : saisie des identifiants, verification BCrypt, emission
d'un Access Token (15 min) et d'un Refresh Token (7 jours), puis
redirection vers le tableau de bord.

\subsection{CU-03 : Voter aux elections}

L'acteur principal est l'Adherent. Les pre-conditions sont : scrutin
en cours, utilisateur inscrit sur les listes electorales \emph{et}
possedant un abonnement saisonnier actif. Le scenario nominal est :
choix d'un candidat, verification d'eligibilite (abonnement actif)
via \code{AuthSubscriptionClient.isActiveSubscriber},
enregistrement du vote, accuse de reception. En cas d'utilisateur
non abonne, le front affiche un bandeau "Abonnement requis pour
voter" avec un CTA "Acheter ma carte".

\subsection{CU-04 : Un parent achete pour son fils (B.18)}

L'acteur principal est le Parent. La pre-condition est que le parent
ait au moins un enfant academie enregistre et un wallet E-Cash
approvisionne. Le scenario nominal est : (1) le parent selectionne
"Pour mon fils" dans le bandeau beneficiaire, (2) le systeme appelle
\code{ensureChildUser} pour creer ou recuperer l'User shadow de
l'enfant, (3) le systeme appelle
\code{SportsAcademyClient.lookup(academyId)} pour verifier que
l'enfant est bien rattache a ce parent (anti-IDOR), (4) le systeme
debite le wallet E-Cash du parent et cree le ticket / abonnement au
nom de l'enfant, avec tracabilite.

\subsection{CU-05 : Convoquer un joueur}

L'acteur principal est l'Entraineur. Le scenario nominal est :
l'entraineur selectionne l'evenement (match ou seance), choisit la
liste des joueurs convoques, definit l'heure, le lieu et la consigne.
Le systeme envoie les convocations (notification in-app et e-mail).

\subsection{CU-06 : Envoyer un message}

L'acteur principal est tout utilisateur authentifie. Le scenario
nominal est : l'utilisateur ouvre ou cree une conversation, saisit son
message et l'envoie. Le message est persiste en base et diffuse en
temps reel via WebSocket STOMP. Une piece jointe peut etre uploadee
au prealable (Cloudinary, URL signee TTL 1h).

\section{Conclusion}

Ce chapitre a permis de formaliser les besoins fonctionnels et non
fonctionnels du systeme, d'identifier les acteurs, et de presenter
les principaux cas d'utilisation. Le chapitre suivant sera consacre a
la conception detaillee du systeme.

% =====================================================================
% CHAPITRE 3
% =====================================================================
\chapterbanner{CHAPITRE 3 : CONCEPTION DU SYSTEME}
\chapter{Conception du système}

\section{Introduction}

Ce chapitre presente la conception detaillee du systeme. On y decrit
l'architecture logicielle et materielle, les principaux diagrammes
(sequence, collaboration, etats, activite, classes), le modele
logique de donnees (MLD) et le dictionnaire des donnees.

\section{Architecture globale}

\begin{keypointbox}
L'architecture de Wydad Digital repose sur \textbf{trois couches
logiques} : (1) \emph{Presentation} avec un frontal Angular 19 (et
trois variantes de layout : public, espace connecte, admin), (2)
\emph{Metier} avec onze microservices~Spring~Boot~3.3 et une passerelle
Spring Cloud Gateway, (3) \emph{Donnees} avec un cluster PostgreSQL
multi-bases, Cloudinary (media) et LiveKit (appels video). Le tout
est conteneurise via Docker Compose et deploye sur une VM Azure
unique.
\end{keypointbox}

\insertimage{archi-microservices.png}{Architecture globale de la
plateforme Wydad Digital. Cette vue presente les acteurs en haut,
le frontal Angular 19 avec 7 espaces utilisateurs (joueur, staff,
entraineur, fan, journaliste, president, admin) + la navigation
publique, la passerelle API Gateway (port 8080) avec authentification
JWT, les onze microservices metier (auth, sports, ticket, payment,
content, communication, election, shop, gamification, notification,
plus la passerelle api-gateway), le cluster PostgreSQL avec 10 bases
logiques, le cache Redis, et les services externes (Cloudinary,
LiveKit). Diagramme PlantUML, sans chevauchement.}

\subsection{Couche Presentation (Frontend Angular 19)}

La couche presentation est entierement implementee en \textbf{Angular
17} (standalone components, signals, route guards, intercepteurs
HTTP). C'est une \emph{single-page application} unique structuree
autour de trois layouts principaux : \code{PublicLayoutComponent},
\code{AdminLayoutComponent}, \code{EspaceLayoutComponent}.

L'organisation du code suit les conventions Angular 19 \emph{standalone}
: \code{app/components/}, \code{app/pages/}, \code{app/layouts/},
\code{app/services/}, \code{app/guards/}, \code{app/interceptors/},
\code{app/directives/}.

Le \textbf{design system} repose sur Tailwind CSS 3.4 avec un jeu de
tokens personnalises (\code{wac-red} = \code{C8102E},
\code{paper-1}, \code{ink-primary}, \code{club-hero}). Le badge dore
GOL \code{D4AF37} est reserve au role president de club.

\subsection{Couche Passerelle (API Gateway)}

L'unique point d'entree externe est l'\textbf{API Gateway} (Spring
Cloud Gateway, Spring Boot 3.3) avec \textbf{JwtAuthenticationFilter},
\textbf{CorsConfig}, blocage des routes \code{/internal/**} depuis
l'exterieur, et health checks Spring Actuator.

\subsection{Couche Metier (11 microservices~Spring~Boot~3.3)}

Chaque microservice Spring Boot 3.3 suit la meme structure de paquets
(\code{controller / service / repository / model / dto / client /
config / filter / util}) et respecte la meme politique de securite.

\begin{table}[H]
\centering
\scriptsize
\caption{Liste des microservices et leur responsabilite}
\label{tab:services}
\begin{tabular}{@{}p{4.3cm}p{2.5cm}p{8.2cm}@{}}
\toprule
\textbf{Service} & \textbf{Base} & \textbf{Role} \\
\midrule
\code{auth-service} (8081) & \code{auth\_db} & Inscription, login, JWT, KYC, abonnements, plans, cartes PDF \\
\code{sports-service} (8087) & \code{sports\_db} & Effectif, convocations, seances, academie \\
\code{ticket-service} (8085) & \code{ticket\_db} & Billetterie, QR codes signes HMAC, scan anti-rejeu \\
\code{payment-service} (8083) & \code{payment\_db} & Paiement, E-Cash, webhooks \\
\code{content-service} (8082) & \code{content\_db} & Actualites, calendrier, partenaires \\
\code{communication-service} (8090) & \code{communication\_db} & Messagerie 1-a-1 et groupe, WebSocket STOMP \\
\code{election-service} (8089) & \code{elections\_db} & Elections, sondages, depouillement \\
\code{shop-service} (8084) & \code{shop\_db} & Catalogue, panier, commandes \\
\code{gamification-service} (8088) & \code{wydad\_gamification} & Points, badges, leaderboard \\
\code{notification-service} (8086) & \code{notification\_db} & Notifications in-app, e-mails \\
\code{api-gateway} (8080) & --- & Routage, securite, CORS \\
\bottomrule
\end{tabular}
\end{table}

\subsection{Couche Donnees (PostgreSQL multi-bases)}

PostgreSQL 16 alpine est deploye dans un conteneur unique avec
\emph{plusieurs bases logiques separees} : \code{auth\_db},
\code{sports\_db}, \code{ticket\_db}, \code{content\_db},
\code{elections\_db}, \code{shop\_db}, \code{payment\_db},
\code{communication\_db}, \code{gamification\_db},
\code{notification\_db}. Chaque microservice ne peut acceder qu'a
\emph{sa} base (isolation des schemas).

\section{Architecture Docker (conteneurisation)}

\insertimage{archi-docker.png}{Architecture Docker du
projet. Cette vue represente : (1) le Dockerfile (recette de build
multi-stage), (2) l'image Docker produite (registre local), (3) le
conteneur instancie (avec volumes, networks, env vars, health
checks), (4) le reseau Docker partage \code{wydad-net}. Le projet
utilise un monorepo avec un \code{docker-compose.yml} unique qui
orchestre 15 conteneurs (11 services metier + 1 api-gateway + 1
frontend + 1 PostgreSQL + 1 Redis).}

\section{Architecture CI/CD (GitHub Actions)}

\insertimage{archi-cicd.png}{Pipeline CI/CD GitHub Actions.
Vue d'ensemble du pipeline en swimlanes : Developpeur, GitHub
Actions, Docker Hub, VM Azure. (1) push sur main, (2) trigger
workflow, (3) jobs paralleles (compilation Maven des 11 services
Spring Boot, build Angular du frontend, push des images Docker sur
le registre local), (4) artefacts uploades, (5) redeployement
manuel sur la VM Azure.}

\section{Architecture Azure (deploiement)}

\insertimage{archi-azure.png}{Infrastructure Azure de Wydad
Digital. Cette vue represente les ressources Azure provisionnees :
(1) VM \code{wydad-vm} (Standard\_B2s\_v2, 2 vCPU, 8~Go RAM, 30~Go
SSD, region Spain Central), (2) Public IP statique
\code{158.158.74.169}, (3) NSG \code{sg-wydad-web} autorisant les
ports 22/TCP (SSH, IP admin), 80/443/TCP (HTTP/HTTPS Internet),
4200/8080/TCP (admin uniquement), (4) Disque SSD Premium 30~Go
(code B2s Premium). Cout total : \textasciitilde{}15 EUR/mois.}

\section{Communication inter-services}

La communication inter-services suit deux modes complementaires :

\begin{enumerate}
  \item \textbf{Synchrone REST} pour les requetes unitaires
        (lecture d'un evenement, verification d'adhesion, lookup
        d'un academicien) via \code{RestTemplate} ou \code{WebClient}.
        Tous les appels internes sont signes par le header
        \code{X-Internal-Secret} (secret partage dans le
        \code{.env}, comparaison en temps constant cote
        microservice recepteur via \code{InternalSecretValidator}).
        La passerelle n'a pas de regle explicite de blocage sur
        \code{/internal/**} : ces routes ne sont pas joignables
        depuis l'exterieur uniquement parce qu'elles n'acceptent
        pas le JWT utilisateur ; la vraie defense en profondeur
        est le \code{X-Internal-Secret} valide cote microservice.
  \item \textbf{Asynchrone WebSocket (STOMP)} pour la messagerie
        temps reel. Le chat 1-a-1 passe par le
        \code{MessagingController} (REST + WebSocket prive),
        tandis que le chat de groupe est gere par le
        \code{TeamChatController} du \code{communication-service}
        (un topic par canal de groupe \code{/topic/team/X})
        avec authentification JWT au CONNECT.
\end{enumerate}

\section{Diagrammes de sequence}

\subsection{Cas nominaux}

\insertimage{sequence-achat-billet-cas1-visiteur.png}{Diagramme de
sequence (Cas 1/4) -- Achat billet par VISITEUR sans compte
(B.28). On voit la creation d'un User guest temporaire, la
reservation du siege, le paiement CMI / PayZone mocke, la
generation du QR code signe, et l'envoi de l'e-mail de confirmation.
Statut billet direct PAID (pas d'en-attente).}

\insertimage{sequence-achat-billet-cas1-paiement.png}{Diagramme
de sequence (Cas 2/4) -- Paiement d'un billet par ADHERENT
identifie. On voit le flow complet de paiement (debit E-Cash
optionnel + paiement carte), la mise a jour du statut ticket
(PAID), la generation du QR code (HMAC-SHA256), et la
notification in-app + e-mail.}

\insertimage{sequence-achat-billet-cas2-parent.png}{Diagramme
de sequence (Cas 3/4) -- Achat PARENT pour LUI-MEME (B.18). Le
parent achete son propre abonnement ou billet, l'API distingue
\texttt{beneficiaryAcademyMemberId=null}.}

\insertimage{sequence-achat-billet-cas2-parent-billet.png}
{Diagramme de sequence (Cas 4/4) -- Achat PARENT pour son
ENFANT (B.18). Le parent selectionne l'enfant dans le bandeau
beneficiaire, l'API resout l'User shadow \code{enfant-
\{academyId\}@wac.parent}, verifie l'anti-IDOR via
\code{SportsAcademyClient.lookup(academyId)} (parent proprietaire
de l'enfant), puis cree le ticket au nom de l'enfant avec
tracabilite \code{parent\_payer\_email}.}

\insertimage{sequence-vote-election.png}{Diagramme de sequence --
Vote a une election presidentielle (cas nominal B.8). Cette vue
presente les acteurs et les lifelines (Frontend, Passerelle,
Service Election, Service Auth, Service Notification). On y voit
la verification de l'adhesion saisonniere ACTIVE en temps reel
(B.8.h, pas fige a startsAt), la verification anti-double-vote
(contrainte SQL UNIQUE sur (election\_id, user\_id)), la
notification in-app de confirmation, et la branche d'erreur 403
VOTE\_REQUIRES\_MEMBERSHIP pour les non-abonnes (CTA <<Acheter
ma carte>>).}

\insertimage{sequence-vote-election-failures.png}{Diagramme de
sequence -- Cas d'echec du vote (B.8). 4 branches : scrutin
FERME (409), scrutin ANNULE, utilisateur non abonne (403
VOTE\_REQUIRES\_MEMBERSHIP), double-vote (409 ALREADY\_VOTED),
auth-service down (fail-closed 503).}

\insertimage{sequence-accreditation-presse.png}{Diagramme de
sequence -- Inscription journaliste et demande d'accreditation
(B.17). Cette vue presente le flux complet : (1) inscription
multipart avec photo + media + n° carte, (2) creation du compte
EN\_ATTENTE, (3) upload photo sur Cloudinary (mode degrade :
secureUrl=null fallback publicId), (4) validation admin (POST
/api/auth/admin/press/accreditations/\{id\}/validate), (5)
generation du badge PDF 19 Ko OpenPDF a la demande (GET
/\{id\}/badge, self ou ADMIN, IDOR strict), (6) telechargement
par journaliste VALIDE.}

\insertimage{sequence-convocations.png}{Diagramme de sequence --
Convoquer un joueur et sa reponse (B.18.b). Cette vue presente le
flux complet : (1) l'entraineur cree une convocation nominative
(coche les joueurs un par un), (2) notification in-app
personnalisee, (3) le joueur accuse reception et repond
(CONFIRME / ABSENT / RETARD avec justification), (4)
l'entraineur consulte le tableau de reponses de sa seance.
Isolation discipline + categorie verifiee cote serveur (IDOR
strict).}

\insertimage{sequence-feuille-match.png}{Diagramme de sequence
-- Feuille de match B.18.b (joueur cocheable). Le capitaine ou
chaque joueur coche sa presence + poste pour le match a venir.
L'entraineur visualise la feuille complete avec statut par
joueur (PRESENT / ABSENT / INCERTAIN). Aggregation temps reel.}

\insertimage{sequence-team-chat.png}{Diagramme de sequence --
Team-chat WebSocket STOMP V2.2. Connexion WS avec JWT au
CONNECT, SUBSCRIBE au topic /topic/team/\{sport\}/\{category\},
envoi via SEND /app/chat/\{sport\}/\{category\}/send, broadcast
aux abonnes. Authentification par sujet : ENTRAINEUR peut
publier uniquement dans son topic discipline+categorie (filter
applique dans le \code{TeamChatAuthInterceptor}).}

\insertimage{sequence-president-receipts.png}{Diagramme de
sequence -- Emission et consultation d'un recu de salaire
(C.21). Le president emet un recu PDF pour un joueur /
entraineur de SA discipline (regle metier anti-fraude), l'API
verifie l'eligibilite, genere le PDF via OpenPDF, le stocke et
renvoie l'URL signee. Le beneficiaire (ou ADMIN) peut
telecharger le PDF ; IDOR strict (proprio ou ADMIN).}

\insertimage{sequence-parent-academie.png}{Diagramme de
sequence -- Inscription d'un enfant a l'academie (B.18). Le
parent remplit le formulaire enfant, l'API cree l'enregistrement
\code{academy\_member} lie a son \code{parentUserId}, et un
User shadow \code{enfant-\{academyId\}@wac.parent} est cree
automatiquement (role VISITEUR + statut VALIDE). Le parent peut
ensuite acheter pour cet enfant.}

\insertimage{sequence-admin-validation-comptes.png}{Diagramme
de sequence -- Validation d'un compte EN\_ATTENTE par ADMIN
(Phase 0). L'admin accede a la liste \code{GET /api/auth/admin/
accounts?statut=EN\_ATTENTE}, valide ou refuse un compte avec
motif. Notification in-app + e-mail envoyes. Statut passe a
VALIDE ou REFUSE. Boucle avec la validation KYC (piece
d'identite).}

\insertimage{sequence-abonnement.png}{Diagramme de sequence --
Achat d'un abonnement saisonnier (B.12). Le flow complet :
selection du plan, debit E-Cash ou paiement carte, creation
\code{UserSubscription} (status ACTIVE, valid\_from/valid\_to),
generation de la carte membre PDF (OpenPDF), 1 seul abonnement
ACTIVE par saison (regle B.12, 409 ALREADY\_SUBSCRIBED en cas
de doublon).}

\subsection{Diagrammes de collaboration}

\insertimage{collab-accreditation.png}{Diagramme de collaboration
-- Inscription journaliste et demande d'accreditation (B.17).
Cette vue met l'accent sur l'ordre temporel des echanges (1.
1...1.N) entre les 4 participants (Journaliste, Frontend,
Passerelle, Auth Service, Notification, Cloudinary) et la base
\code{auth\_db}. Format Visual Paradigm, focus sur les
messages.}

\section{Diagrammes d'etats (cycles de vie)}

\begin{keypointbox}
Les diagrammes suivants representent les \textbf{cycles de vie
metier} tels qu'ils sont implementes dans le code qui compile
(etat des enums JPA, methodes \code{Service} de transition,
garde-fous et exceptions 409/403). Chaque transition correspond
a un endpoint REST ou un evenement planifie (\code{@Scheduled}).
\end{keypointbox}

\insertimage{state-ticket.png}{Diagramme d'etats (1/3) -- Cycle
de vie d'un Billet. 5 etats : RESERVE (apres creation), PAYE
(apres paiement valide), VALIDE (apres scan QR a l'entree),
ANNULE (annulation avant ou apres paiement). L'etat final est
consomme par scan QR au stade.}

\insertimage{state-election.png}{Diagramme d'etats (2/3) --
Cycle de vie d'une Election (B.8 / B.8.b / B.8.c / B.8.e /
B.8.h). 3 etats : OPEN (published=false, votes acceptes),
CLOSED (published=false, gele, pas de resultats), PUBLISHED
(status=CLOSED, published=true, winnerCandidateId fige).
\textbf{Regles B.8} : \emph{close} gele SANS publier
(\code{ElectionService.closeOnlyInternal}), \emph{publish}
exige TOUS titulaires votes (409 NOT\_ALL\_VOTED) et
eligibles > 0 (409 PUBLISH\_NO\_ELIGIBLE\_SNAPSHOT, fail-closed
si auth-service down). \textbf{B.8.b} \emph{unpublish}
ADMIN fait repasser published=false (status reste CLOSED,
winnerCandidateId conserve). \textbf{B.8.c} garde-fou
$\geq 2$ candidats avant close. \textbf{B.8.e} \code{toggleResultsHidden}
interdit apres publication. \textbf{B.8.h} \code{canVote}
= titulaire ACTIF en temps reel (pas fige a startsAt) + ADMIN
exempte, 1 vote par (election, userId). Scheduler
\code{@Scheduled fixedDelay=60s}.}

\insertimage{state-press-accred.png}{Diagramme d'etats (3/3) --
Cycle de vie d'une PressAccreditation (B.17). 3 etats :
EN\_ATTENTE (journaliste cree, multipart photo + media + n°
carte, statut EN\_ATTENTE, dateCreation=now), VALIDE
(ADMIN /validate, decidedAt=now, motifRefus=null), REFUSE
(ADMIN /refuse, motifRefus obligatoire \code{@Valid size>=5},
decidedAt=now). Apres validation : badge PDF 19 Ko
telechargeable (GET /accreditations/\{id\}/badge, OpenPDF,
self OU ADMIN, 403 IDOR strict sinon). Apres refus : motif
visible cote journaliste, notif in-app, \textbf{PAS de
re-soumission} sur la meme demande (POST /accreditations
pour repartir d'EN\_ATTENTE).}

\section{Diagramme de classes (vue macro)}

\insertimage{classes-macro.png}{Diagramme de classes -- Vue
macro du domaine Wydad Digital. Cette vue presente les principales
classes metier regroupees par domaine : Auth (User, Subscription,
PressAccreditation), Billetterie (Event, Section, Ticket),
Boutique (Product, Order, OrderItem), Sports (Player, Convocation,
Session), Elections (Election, Candidate, Vote). Format Visual
Paradigm avec packages, attributs essentiels, et relations
(association, agregation, composition) avec cardinalites.}

\section{Diagrammes de classes par domaine}

\insertimage{classes-auth.png}{Diagramme de classes (Domaine
1/6) -- Auth \& Identite. Classes : \code{User},
\code{ActiveSession}, \code{KycDocument},
\code{SubscriptionPlan}, \code{UserSubscription},
\code{PressAccreditation}, \code{SalaryReceipt}, avec leurs
attributs, enums associes (Role, StatutCompte, MembershipLevel,
etc.) et les relations JPA (composition, agregation, FK).}

\insertimage{classes-billetterie.png}{Diagramme de classes
(Domaine 2/6) -- Billetterie \& Contenu. Classes : \code{Event},
\code{Section}, \code{Ticket}, \code{Match}, \code{Article}, avec
leurs enums (EventStatus, TicketStatus, TicketCategory, etc.)
et les relations JPA (FK logique vers Match content-service).}

\insertimage{classes-sports.png}{Diagramme de classes
(Domaine 3/6) -- Sports \& Roster. Classes : \code{Player},
\code{Staff}, \code{AcademyMember}, \code{Session},
\code{Convocation}, \code{MatchConvocation}, \code{ScheduledCall},
avec leurs enums (SportType, Category, MedicalStatus, etc.) et
les relations JPA.}

\insertimage{classes-shop.png}{Diagramme de classes (Domaine
4/6) -- Boutique \& Paiement. Classes : \code{Category},
\code{Product}, \code{ProductVariant}, \code{ProductImage},
\code{ProductReview}, \code{Store}, \code{PromoCode},
\code{JerseyCustomization}, \code{ShopOrder}, \code{OrderItem},
\code{OrderStatusHistory}, \code{ECashAccount}, \code{Transaction},
avec leurs enums et les relations JPA (composition ALL +
orphanRemoval).}

\insertimage{classes-academie.png}{Diagramme de classes
(Domaine 5/6) -- Academie (B.18). Classes : \code{AcademyMember},
\code{User} shadow \code{enfant-\{academyId\}@wac.parent},
\code{SportsAcademyClient} (lookup anti-IDOR), \code{AcademyController}
(B.18 : CRUD admin academie), \code{MatchConvocation}
(convocations nominatives B.18.b).}

\insertimage{classes-messagerie.png}{Diagramme de classes
(Domaine 6/6) -- Messagerie V2.5. Classes : \code{Message}
(1-to-1 persistant avec photo), \code{TeamMessage} (groupe
STOMP), \code{Attachment} (Cloudinary publicId + secureUrl),
\code{MessageMediaService} (V2.3 upload), \code{TeamChatService}
(V2.5 sendWithMedia), \code{MessagingService} (1-to-1 IDOR
strict sameGroup), \code{RosterClient} + \code{AuthUserInfoClient}
(resolution president discipline).}

\insertimage{classes-election.png}{Diagramme de classes
(Domaine transverse) -- Elections B.8. Classes : \code{Election}
(status + published + resultsHidden + winnerCandidateId),
\code{Candidate}, \code{Vote} (contrainte UNIQUE
(election\_id, user\_id)), \code{ActiveMembersClient}
(snapshot eligibilite B.8.h), \code{ElectionService}
(closeOnlyInternal, publishResultsInternal, unpublishForAdmin,
toggleResultsHidden, closeAndPublish).}

\insertimage{classes-gamification.png}{Diagramme de classes
(Domaine transverse) -- Gamification C1/C2/C3. Classes :
\code{PointTransaction}, \code{UserPoints}, \code{PointRule}
(C3 bareme, cache memoire), \code{Challenge} (C2 defis
temporels), \code{UserChallengeProgress}, \code{Badge},
\code{UserBadge}, \code{Pronostic} (C1), \code{GamificationService}
(attribution automatique sur evenements).}

\section{Diagramme de classes (vue macro)}

\insertimage{classes-macro.png}{Diagramme de classes -- Vue
macro du domaine Wydad Digital. Cette vue presente les principales
classes metier regroupees par domaine : Auth (User, Subscription,
PressAccreditation), Billetterie (Event, Section, Ticket),
Boutique (Product, Order, OrderItem), Sports (Player, Convocation,
Session), Elections (Election, Candidate, Vote). Format Visual
Paradigm avec packages, attributs essentiels, et relations
(association, agregation, composition) avec cardinalites.}

\section{Diagrammes de classes par domaine}

\insertimage{classes-auth.png}{Diagramme de classes -- Domaine
Auth \& Identite. Cette vue detaillee presente les classes :
\code{User}, \code{ActiveSession}, \code{KycDocument},
\code{SubscriptionPlan}, \code{UserSubscription},
\code{PressAccreditation}, \code{SalaryReceipt}, avec leurs
attributs, enums associes (Role, StatutCompte, MembershipLevel,
etc.) et les relations JPA (composition, agregation, FK).}

\insertimage{classes-billetterie.png}{Diagramme de classes --
Domaine Billetterie \& Contenu. Cette vue detaillee presente les
classes : \code{Event}, \code{Section}, \code{Ticket},
\code{Match}, \code{Article}, avec leurs enums (EventStatus,
TicketStatus, TicketCategory, etc.) et les relations JPA (FK
logique vers Match content-service).}

\insertimage{classes-sports.png}{Diagramme de classes -- Domaine
Sports \& Roster. Cette vue detaillee presente les classes :
\code{Player}, \code{Staff}, \code{AcademyMember}, \code{Session},
\code{Convocation}, \code{MatchConvocation}, \code{ScheduledCall},
avec leurs enums (SportType, Category, MedicalStatus, etc.) et les
relations JPA.}

\insertimage{classes-shop.png}{Diagramme de classes -- Domaine
Boutique \& Paiement. Cette vue detaillee presente les classes :
\code{Category}, \code{Product}, \code{ProductVariant},
\code{ProductImage}, \code{ProductReview}, \code{Store},
\code{PromoCode}, \code{JerseyCustomization}, \code{ShopOrder},
\code{OrderItem}, \code{OrderStatusHistory}, \code{ECashAccount},
\code{Transaction}, avec leurs enums et les relations JPA (composition
ALL + orphanRemoval).}

\section{Modele logique de donnees (MLD)}

\begin{table}[H]
\centering
\scriptsize
\caption{Modele logique de donnees (extrait representatif)}
\label{tab:mld}
\begin{tabular}{@{}p{5.5cm}p{9.5cm}@{}}
\toprule
\textbf{Table} & \textbf{Description} \\
\midrule
\code{users} & Comptes utilisateurs (id, email, password, role, statut\_compte, firstName, lastName, phone, discipline\_demandee, categorie\_demandee) \\
\code{user\_subscriptions} & Abonnements saisonniers (id, user\_id, plan\_id, season, paid\_amount, paid\_at, valid\_from, valid\_to, status, beneficiary\_academy\_member\_id, parent\_payer\_email) \\
\code{tickets} & Billets achetes (id, ticket\_number, user\_id, event\_id, section\_id, qr\_code, status, beneficiary\_academy\_member\_id, parent\_payer\_email) \\
\code{events} & Matchs / evenements (id, title, event\_date, venue, description, available\_seats, sold\_tickets) \\
\code{sections} & Sections de tribune (id, event\_id, name, price, available\_seats) \\
\code{academy\_members} & Enfants inscrits a l'academie (id, parent\_user\_id, child\_full\_name, child\_birth\_date, sport\_type, level, active) \\
\code{players} & Fiches roster joueurs (id, user\_id, full\_name, sport\_type, category, position, jersey\_number, photo\_url) \\
\code{staff} & Fiches roster staff (id, user\_id, full\_name, sport\_type, assigned\_category, role) \\
\code{elections} & Scrutins (id, title, starts\_at, ends\_at, status, published, winner\_candidate\_id) \\
\code{election\_candidates} & Candidats (id, election\_id, full\_name, presentation, photo\_url, display\_order) \\
\code{election\_votes} & Votes (election\_id, user\_id, candidate\_id) --- contrainte UNIQUE (election\_id, user\_id) \\
\code{polls} / \code{poll\_options} / \code{poll\_votes} & Sondages \\
\code{ecash\_accounts} / \code{ecash\_transactions} & Portefeuille interne E-Cash (id, email, balance, active) \\
\bottomrule
\end{tabular}
\end{table}

\section{Dictionnaire des donnees}

\begin{landscape}
\footnotesize
\begin{longtable}{@{}>{\columncolor{wacGrisClair}}p{4.5cm}p{3.5cm}p{4cm}p{11cm}@{}}
\hline
\rowcolor{wacRouge}
\textcolor{wacBlanc}{\textbf{Champ}} & \textcolor{wacBlanc}{\textbf{Type}} & \textcolor{wacBlanc}{\textbf{Contraintes}} & \textcolor{wacBlanc}{\textbf{Description}} \\ \hline
\endhead
\code{id} & BIGSERIAL & PK & Identifiant unique \\ \hline
\code{email} & VARCHAR(255) & UNIQUE, NOT NULL & Adresse e-mail (login) \\ \hline
\code{password} & VARCHAR(255) & NOT NULL & Mot de passe hache BCrypt \\ \hline
\code{role} & VARCHAR(20) & NOT NULL, CHECK & Role utilisateur (7 valeurs + ADMIN) \\ \hline
\code{statut\_compte} & VARCHAR(16) & NOT NULL, CHECK & EN\_ATTENTE / VALIDE / REFUSE \\ \hline
\code{firstName} & VARCHAR(80) & NOT NULL & Prenom \\ \hline
\code{lastName} & VARCHAR(80) & NOT NULL & Nom \\ \hline
\code{phone} & VARCHAR(32) & UNIQUE & Telephone \\ \hline
\code{discipline\_demandee} & VARCHAR(32) & NULLABLE & Discipline (FOOTBALL, etc.) \\ \hline
\code{categorie\_demandee} & VARCHAR(16) & NULLABLE & Categorie (U15, U17, U18, U20, SENIOR) \\ \hline
\code{beneficiary\_academy\_member\_id} & BIGINT & NULLABLE & B.18 : id de l'enfant (achat pour fils) \\ \hline
\code{parent\_payer\_email} & VARCHAR(256) & NULLABLE & B.18 : email parent payeur \\ \hline
\code{createdAt} & TIMESTAMP & NOT NULL & Date de creation \\ \hline
\code{updatedAt} & TIMESTAMP & NOT NULL & Date de mise a jour \\ \hline
\end{longtable}
\end{landscape}

\section{Conclusion}

Ce chapitre a presente la conception detaillee du systeme, en
couvrant l'architecture logicielle et materielle, les principaux
diagrammes UML, ainsi que le modele de donnees. Le chapitre suivant
sera consacre a la realisation detaillee par couche technique.

% =====================================================================
% CHAPITRE 4
% =====================================================================
\chapterbanner{CHAPITRE 4 : REALISATION DU SYSTEME}
\chapter{Réalisation du système}

\section{Introduction}

Ce chapitre presente la phase de realisation du systeme, structuree
par couche technique : (1) backend Spring Boot, (2) frontend
Angular, (3) conteneurisation Docker, (4) pipeline CI/CD, (5)
deploiement Azure. Pour chaque couche, on decrit l'environnement, les
fichiers de configuration cles, et les principales interfaces
utilisateur.

\section{Couche 1 : Backend Spring Boot 3.3}

\subsection{Environnement backend}

L'environnement backend est constitue de Java OpenJDK 21 LTS (Corretto)
en build Maven 3.9.12 (monorepo multi-modules). Le framework
principal est Spring Boot 3.3.5, accompagne de Spring Cloud
2023.0.3, Spring Security 6.3, et Spring Data JPA 3.3. L'ORM utilise
est Hibernate 6.5 avec Lombok 1.18.30. La base de donnees est
PostgreSQL 16 (alpine). La securite repose sur JWT (jjwt 0.12.x),
BCrypt (cout 12), et HMAC-SHA256 pour les QR codes. L'IDE utilise
est IntelliJ IDEA 2025.2.

\subsection{Structure du monorepo Maven}

L'ensemble des 11 microservices Spring Boot est organise en monorepo
Maven, avec un POM parent \code{wydad-digital-parent} qui
herite de Spring Boot Starter Parent et declare les versions communes.
Chaque microservice possede son propre sous-module Maven avec son
\code{pom.xml}, ses sources \code{src/main/java} et
\code{src/test/java}, et son \code{Dockerfile}.

\subsection{Pattern de securite uniforme}

Chaque microservice applique le meme pattern de securite :

\begin{itemize}
  \item \textbf{\code{@PreAuthorize}} sur chaque methode
        \code{controller} pour restreindre l'acces par role.
  \item \textbf{\code{UserContextFilter}} qui lit les en-tetes
        \code{X-User-*} propages par la passerelle et les
        stocke dans un \code{ThreadLocal} accessible via
        \code{UserContext.getCurrentUserId()}, etc.
  \item \textbf{Validation cote serveur systematique} avec Jakarta
        Validation.
  \item \textbf{\code{GlobalExceptionHandler}} qui centralise la
        transformation des exceptions metier en codes HTTP
        (403 VOTE\_REQUIRES\_MEMBERSHIP, 409 ALREADY\_SUBSCRIBED,
        404 NOT\_YOUR\_CHILD, etc.).
  \item \textbf{\code{SecurityConfig}} qui configure
        \code{csrf().disable()},
        \code{sessionCreationPolicy(STATELESS)} et
        \code{addFilterBefore(userContextFilter, ...)}.
\end{itemize}

\subsection{Endpoints internes inter-services}

\begin{table}[H]
\centering
\scriptsize
\caption{Principaux endpoints internes inter-services}
\label{tab:internals}
\begin{tabular}{@{}p{6.8cm}p{4.5cm}p{3.7cm}@{}}
\toprule
\textbf{Endpoint} & \textbf{Service appelant} & \textbf{But} \\
\midrule
\code{GET /api/auth/internal/is-active-subscriber} & election-service & B.18 : eligibilite vote \\
\code{GET /api/auth/internal/users/\{id\}/discipline} & sports-service, comm-service & C.21 : role + discipline \\
\code{POST /api/auth/internal/ensure-child-user} & ticket-service & B.18 : User shadow enfant \\
\code{GET /api/sports/academy/internal/\{id\}} & auth-service, ticket-service & B.18 : lookup academie \\
\code{POST /api/payment/internal/credit} & scripts de setup & Crediter wallet E-Cash \\
\code{GET /api/sports/internal/roster/membership/\{id\}} & communication-service & Resoudre fiche roster \\
\bottomrule
\end{tabular}
\end{table}

Tous ces endpoints ne sont \textbf{pas joignables depuis un
navigateur} (la passerelle ne valide pas de JWT utilisateur sur
ces chemins) et ne sont accessibles qu'entre microservices
via le reseau Docker interne, avec le header
\code{X-Internal-Secret} verifie en temps constant cote
recepteur par \code{InternalSecretValidator} (defense en
profondeur). Le secret est charge depuis
\code{\$\{INTERNAL\_SECRET\}} du \code{.env}.

\subsection{Securite des mots de passe et JWT}

Les mots de passe sont hashes par BCrypt avec un cout de 12. Les
tokens JWT (Access Token, duree 15 min ; Refresh Token, duree 7
jours) sont signes en HS256 avec le secret \code{JWT\_SECRET}
(256 bits minimum). Le \code{JwtAuthenticationFilter} de la
passerelle supprime systematiquement les en-tetes \code{X-User-*}
fournis par le client (anti-spoofing) et les reconstruit a partir
des claims JWT valides.

\section{Couche 2 : Frontend Angular 19}

\subsection{Environnement frontend}

L'environnement frontend utilise Node.js 20.18 LTS avec npm 10.9 et
Angular CLI 17.0.0. Le framework est Angular 19 (standalone
components, signals, route guards, intercepteurs HTTP). L'UI repose
sur Tailwind CSS 3.4 avec des tokens personnalises. RxJS 7.8 est
utilise pour la programmation reactive.

\subsection{Architecture du frontend}

L'application Angular 19 est organisee en standalone components. La
structure de dossiers : \code{app/components/}, \code{app/pages/},
\code{app/layouts/}, \code{app/services/}, \code{app/guards/},
\code{app/interceptors/}.

\subsection{Securite cote frontend}

L'\code{authInterceptor} ajoute automatiquement le header
\code{Authorization: Bearer ...} a chaque requete sortante. L'
\code{errorInterceptor} centralise la gestion des 401 (refresh
token automatique), 403 et 500. Neuf gardes de route protegent
les espaces.

\section{Couche 3 : Conteneurisation Docker}

L'application est entierement conteneurisee via \textbf{Docker
Compose}. Chaque service dispose de son \code{Dockerfile} (build
multi-stage pour le frontend Angular, image
\texttt{eclipse-temurin:21-jre} pour les Spring Boot). Les
conteneurs partagent un reseau Docker dedie \code{wac-net}, un
volume nomme \code{pgdata}, et des health checks Spring Actuator
avec politique de redemarrage \code{unless-stopped}.

\section{Couche 4 : Pipeline CI/CD (GitHub Actions)}

Le pipeline CI/CD est implemente via \textbf{GitHub Actions} dans
\code{.github/workflows/ci.yml}. A chaque \code{push} sur la
branche \code{main}, le pipeline execute en sequence : (1)
build parent POM, (2) compilation Maven des 11 services Spring
Boot, (3) build + push des images Docker, (4) build du frontend
Angular.
Le redeploiement sur la VM Azure est manuel : \code{ssh
azureuser@158.158.74.169 \&\& git pull \&\& mvn -q -T 1C
package \&\& docker compose up -d --force-recreate
--no-deps --build}.

\section{Couche 5 : Deploiement Azure}

L'infrastructure de production est une \textbf{VM Azure unique} de
gamme \textbf{Standard\_B2s\_v2} (2 vCPU, 8~Go RAM, 30~Go SSD
Premium), situee en region \textbf{Spain Central}, avec une
adresse IP publique statique \code{158.158.74.169}. Le cout
mensuel est de l'ordre de 30~\$. Le systeme d'exploitation est
\textbf{Ubuntu 22.04 LTS Server}. Docker Engine 24.x et Docker
Compose v2.x sont installes. Les sauvegardes quotidiennes
\code{pg\_dumpall} sont effectuees via cron a 03h00 et
poussees vers un compte de stockage Azure Blob.

\section{Principales interfaces graphiques}

\begin{infobox}
Cette section enumere les \textbf{interfaces reellement
implementees} dans le code source du projet (fichiers
\code{wydad-frontend/src/app/pages/} et
\code{wydad-frontend/src/app/components/}). Chaque capture
d'ecran est identifiee par le \textbf{nom de son fichier PNG}, a
placer a la racine du projet Overleaf (au meme niveau que
\code{main.tex}) pour remplacer la cadre gris par la capture reelle.
Les routes Angular correspondantes sont precisees pour permettre
la navigation manuelle lors de la prise de vue.
\end{infobox}

\subsection{Pages publiques (non authentifie)}

Le \code{HomeComponent} (route \code{/}) sert de page d'accueil et
regroupe en un seul ecran : hero classement, prochain match, palmares,
partenaires officiels et carousel d'actualites.

\insertimage{accueil.png}{\code{HomeComponent} (route \code{/}).
Page d'accueil publique du site du WAC : hero classement, prochain
match, palmares, partenaires officiels, carousel d'actualites.}

\insertimage{inscription.png}{\code{RegisterComponent} (route
\code{/register}). Formulaire d'inscription multi-statuts avec
selection du role (VISITEUR, ADHERENT, PARENT, JOUEUR, ENTRAINEUR,
STAFF, JOURNALISTE, PRESIDENT), depot d'un justificatif d'identite
pour les roles privilegies, et selection de la discipline / categorie.}

\insertimage{connexion.png}{\code{LoginComponent} (route
\code{/login}). Page de connexion email + mot de passe avec lien
<< mot de passe oublie >> et CTA << creer un compte >>.}

\insertimage{actualites.png}{\code{ActualitesComponent} (route
\code{/actualites}). Liste paginee des actualites du club, avec
filtre par categorie (match, transfer, club, academie).}

\insertimage{actualite_detail.png}{\code{ActualiteDetailComponent}
(route \code{/actualites/:id}). Detail d'une actualite avec image
hero, contenu enrichi, et liste des articles lies.}

\insertimage{matchs.png}{\code{MatchesComponent} (route
\code{/matchs}). Calendrier des matchs avec filtres (saison,
competition, equipe), score final et billetterie rapide.}

\insertimage{classements.png}{\code{ClassementComponent} (route
\code{/classements}). Classement public de la Botola Pro avec
points, victoires, nuls, defaites et difference de buts.}

\insertimage{billetterie.png}{\code{BilletterieComponent} (route
\code{/billetterie}). Catalogue billetterie : evenements ouverts
avec date, lieu, equipe adverse, plan de salle et nombre de places
disponibles par section.}

\insertimage{billetterie_detail.png}{\code{BilletterieDetailComponent}
(route \code{/billetterie/:id}). Detail d'un evenement avec
selection de section, nombre de places, recapitulatif et CTA
<< Acheter >> (bandeau parent B.18 si applicable).}

\insertimage{abonnement.png}{\code{AbonnementComponent} (route
\code{/abonnement}). Catalogue des plans d'abonnement saisonnier
avec prix regulier, prix adherent et bandeau B.18
<< Pour moi / Pour mon fils >> visible pour les parents.}

\insertimage{legendes.png}{\code{LegendesComponent} (route
\code{/legendes}). Vitrine des legendes du WAC avec photo,
periode, palmares individuel et anecdotes.}

\insertimage{palmares.png}{\code{PalmaresComponent} (route
\code{/palmares}). Palmares detaille du club (Botola, Coupe du
Trone, Ligue des Champions CAF) par saison et competition.}

\insertimage{stade.png}{\code{StadeComponent} (route \code{/stade}).
Presentation du stade (capacite, plan, histoire).}

\insertimage{transparence.png}{\code{TransparenceComponent} (route
\code{/transparence}). Page publique transparence (rapports
financiers, gouvernance).}

\subsection{Boutique (publique + authentifie)}

\insertimage{catalogue_produit.png}{\code{BoutiqueComponent} (route
\code{/boutique}). Catalogue produits : filtres par categorie,
sport, recherche textuelle, image / nom / prix / bouton
<< Ajouter au panier >>.}

\insertimage{boutique_detail.png}{\code{BoutiqueDetailComponent}
(route \code{/boutique/:id}). Fiche produit avec galerie images,
variantes (taille / couleur), personnalisation maillot, avis
clients et bouton << Ajouter au panier >>.}

\insertimage{panier.png}{\code{CartComponent} (route \code{/panier},
auth requise). Panier d'achat : liste des articles, quantite, prix
unitaire, prix total, boutons << Vider le panier >> et
<< Passer commande >>.}

\insertimage{historique_commandes.png}{\code{MesCommandesComponent}
(route \code{/profil/commandes}, auth requise). Historique des
commandes (EN\_ATTENTE / PAYEE / EXPEDIEE / LIVREE) avec date,
montant, telechargement de la facture PDF.}

\insertimage{mes_achats.png}{\code{MesAchatsComponent} (route
\code{/profil/ecash}, auth requise). Historique unifie des achats
(billets, abonnements, dons, boutique) avec filtre par type et
date.}

\subsection{Espace adherent / supporter}

\insertimage{espace_fan.png}{\code{EspaceFanComponent} (route
\code{/espace-fan}, auth requise). 4 onglets : \textbf{Overview}
(classement personnel, prochain match, prochain defi, prochaines
predictions), \textbf{Predictions} (pronostics sur matchs a venir),
\textbf{Defis} (defis C2 en cours, progression), \textbf{Gains}
(historique des 50 dernieres transactions de points, solde
personnel).}

\insertimage{mes_billets.png}{\code{MesBilletsComponent} (route
\code{/profil/billets}, auth requise). Billets achetes avec QR code,
statut (PAID / USED / CANCELLED), date, prix, bouton
<< Telecharger PDF >>. Un parent y retrouve aussi les billets
achetes au nom de ses enfants (B.18).}

\insertimage{ma_carte.png}{\code{CarteMembreComponent} (route
\code{/profil/carte}, auth requise). Carte de membre numerique avec
QR code d'identification, numero d'adherent, plan et date
d'expiration.}

\insertimage{mon_ecash.png}{\code{EcashComponent} (route
\code{/profil/ecash}, auth requise). Portefeuille E-Cash : solde,
historique des transactions, bouton << Approvisionner >>.}

\insertimage{mon_profil.png}{\code{ProfilComponent} (route
\code{/profil}, auth requise). Edition des informations
personnelles, photo de profil, preferences de notification,
changement de mot de passe.}

\insertimage{mon_don.png}{\code{DonComponent} (route \code{/don},
auth requise). Page de don au club avec saisie du montant, choix
du mode de paiement (E-Cash / carte) et confirmation.}

\insertimage{mes_elections.png}{\code{MesElectionsComponent} (route
\code{/mes-elections}, auth requise). Page << Mes elections >> avec
bandeau 403 VOTE\_REQUIRES\_MEMBERSHIP si l'utilisateur n'a pas
d'abonnement saisonnier actif (CTA << Acheter ma carte >>).}

\insertimage{elections_liste.png}{\code{ElectionsComponent} (route
\code{/elections}). Liste publique des elections en cours /
passees avec dates, statut (OUVERTE / FERMEE) et resultats si
disponibles.}

\insertimage{vote_election.png}{Interface de vote dans
\code{MesElectionsComponent} : liste des candidats (photo +
presentation), selection d'un candidat, verification d'adhesion
saisonniere cote back, enregistrement et accuse de reception.}

\insertimage{resultats_election.png}{Page publique des resultats
d'une election (dans \code{ElectionsComponent}) : diagramme en
barres des voix par candidat, pourcentages, gagnant declare.
Accessible sans authentification.}

\insertimage{liste_sondages.png}{\code{SondagesComponent} (route
\code{/sondages}). Liste publique des sondages actifs : question,
options, date de fermeture.}

\insertimage{vote_sondage.png}{Interface de vote a un sondage a
choix multiples dans \code{SondagesComponent} : selection d'une
option, un seul vote par utilisateur, accuse de reception.}

\subsection{Espace parent (B.18)}

\insertimage{espace_parent.png}{\code{DashboardParentComponent}
(route \code{/academie/mes-enfants}, auth PARENT). Section
<< Espace Supporter >> en haut avec 5 tuiles (Abonnement,
Billetterie, Mes billets, Ma carte, Elections). En dessous, la liste
des enfants inscrits a l'academie avec planning d'entrainement
cliquable, statut ACTIF / EN\_ATTENTE, et bouton
<< Telecharger pieces >>.}

\insertimage{inscription_enfant.png}{\code{InscriptionAcademieComponent}
(route \code{/academie/inscription}, auth PARENT ou ADMIN).
Formulaire d'inscription d'un enfant : informations personnelles,
selection de la discipline et de la categorie, depot des pieces
justificatives (acte de naissance, certificat medical, photo).}

\subsection{Espace joueur}

\insertimage{espace_joueur.png}{\code{DashboardJoueurComponent}
(route \code{/joueur/dashboard}, auth JOUEUR). Header pro
(breadcrumb + identite + badge medical), KPI 4 colonnes, sidebar
raccourcis, 4 onglets principaux : \textbf{Quotidien} (seances
du jour, prochaines convocations, matchs a venir, resume mes
billets VIP), \textbf{Convos} (table des convocations match et
seances avec reponse CONFIRME / ABSENT / RETARD), \textbf{Mes
docs} (documents KYC deposes, statut), \textbf{Comms} avec
sous-onglets V2.5 (\textbf{Mon groupe} = team-chat enrichi photos,
\textbf{Canal president} = 1-to-1 persistant, bouton
\textbf{Nouveau message} = selecteur libre).}

\subsection{Espace staff}

\insertimage{espace_staff.png}{\code{DashboardStaffComponent}
(route \code{/staff/dashboard}, auth STAFF). Sidebar 4 onglets :
\textbf{Quotidien} (resume effectif, seances, convocations, KPI),
\textbf{Comms} (memes sous-onglets V2.5 que joueur :
Mon groupe / Canal president / Nouveau message), \textbf{Media}
(mediateque tactique, upload video / schema / note), \textbf{Appels}
(programmation d'appels video LiveKit, agenda partage).}

\subsection{Espace entraineur}

\insertimage{espace_entraineur.png}{\code{DashboardEntraineurComponent}
(route \code{/entraineur/dashboard}, auth ENTRAINEUR). 7 onglets
metier : \textbf{Effectif} (liste joueurs de SA categorie avec
position, jersey number, fiche), \textbf{Seances} (creation et suivi
des seances d'entrainement), \textbf{Matchs} (convocations match
batch, feuille de match), \textbf{Messagerie} (V2.4 batch 1..N avec
photo), \textbf{Reunions video} (programmation LiveKit),
\textbf{Recus} (C.21 recus salaire), \textbf{VIP} (vue rapide des
billets VIP emis). L'onglet Messagerie a 3 sous-onglets
\textbf{groupe} / \textbf{entraineur} / \textbf{president} +
bouton \textbf{Nouveau message}.}

\insertimage{convocation_seance.png}{Formulaire de creation d'une
seance d'entrainement dans \code{DashboardEntraineurComponent} :
date, heure, lieu, type (entrainement / match), liste des joueurs
convoques, consigne. Gestion des reponses (CONFIRME / ABSENT /
RETARD avec justification obligatoire).}

\insertimage{convocation_match.png}{Formulaire de convocation
d'un match dans \code{DashboardEntraineurComponent} : selection du
match, heure, lieu, consigne tactique, selection multiple des
joueurs (batch).}

\insertimage{fiche_joueur.png}{Fiche individuelle d'un joueur dans
l'onglet \textbf{Effectif} : poste, numero de maillot, taille,
pointure, poids, observations medicales, statistiques (buts,
passes decisives, matches joues), historique des convocations.}

\insertimage{feuille_match.png}{Feuille de match electronique dans
l'onglet \textbf{Matchs} : emargement (presents / absents /
excuses), selection du poste de chaque joueur, validation
ADMIN avant publication.}

\insertimage{mediateque.png}{Mediateque tactique d'une equipe dans
l'onglet \textbf{Media} : liste des documents partages (videos,
schemas, notes), upload, partage avec un joueur specifique ou toute
l'equipe.}

\subsection{Espace journaliste (B.17)}

\insertimage{espace_journaliste.png}{\code{DashboardJournalisteComponent}
(route \code{/journaliste/accueil}, auth JOURNALISTE). 4 onglets :
\textbf{Badge} (badge PDF officiel telechargeable, signature
serveur), \textbf{Matchs} (matchs couverts, planning), \textbf{Demandes}
(historique des demandes d'accreditation avec statut), \textbf{Actualites}.}

\insertimage{mes_demandes_journaliste.png}{\code{MesDemandesJournalisteComponent}
(route \code{/journaliste/demandes}, auth JOURNALISTE). Historique
complet des demandes d'accreditation avec statut (EN\_ATTENTE /
VALIDE / REFUSE), motif de refus si applicable, badge PDF si
valide.}

\insertimage{demande_accreditation.png}{Formulaire de demande
d'accreditation (dans \code{DashboardJournalisteComponent} onglet
\textbf{Demandes}) : informations sur le media, numero de carte
de presse, photo, piece jointe.}

\insertimage{badge_presse.png}{Badge d'accreditation presse genere
en PDF (Telechargeable depuis l'onglet \textbf{Badge}) : nom du
journaliste, organe de presse, n° de carte de presse, match couvert,
QR code de controle. Le badge est signe numeriquement cote
serveur.}

\subsection{Espace president de club (C.21)}

\insertimage{dashboard_president.png}{\code{PresidentDashboardComponent}
(route \code{/president/dashboard}, auth PRESIDENT). Bandeau
\emph{dore GOL} (au lieu du rouge club standard) avec badge
<< GOL >>, nom du president, discipline, et regle d'or :
<< Vous ne communiquerez qu'avec les joueurs et entraineurs de
cette discipline >>. En dessous, selecteur discipline / categorie
+ 4 onglets : \textbf{Messagerie} (team-chat + 1-to-1 + selecteur
libre V2.5), \textbf{Annuaire} (joueurs + entraineurs de SA
discipline), \textbf{Recus} (C.21 emission PDF recus salaires),
\textbf{Video} (programmation appels video LiveKit).}

\insertimage{president_discussions.png}{\code{PresidentDiscussionsComponent}
(route \code{/president/discussions}, auth PRESIDENT). Vue
discussion dediee avec liste des conversations 1-to-1 et bouton
<< Message libre >> qui ouvre le selecteur libre V2.5
(\code{MessagerieLibreComponent}) pre-rempli avec le destinataire
selectionne.}

\insertimage{annuaire_president.png}{Onglet \textbf{Annuaire} de
\code{PresidentDashboardComponent} : liste filtree des joueurs et
entraineurs de SA discipline (FOOTBALL / BASKETBALL / HANDBALL)
uniquement. Chaque tuile propose un bouton
<< Message >> (bascule vers l'onglet Messagerie avec destinataire
pre-rempli) et un bouton << Appeler >> (pre-remplit
\code{ScheduleCallFormComponent} avec \code{initialTargetUserIds=[id]}).}

\insertimage{recu_salaire.png}{Onglet \textbf{Recus} de
\code{PresidentDashboardComponent} : emission d'un recu de salaire
ou de prime. Selection du beneficiaire (joueur ou staff de SA
discipline), montant, periode, motif. Le recu est genere en PDF
(OpenPDF) et telechargeable.}

\insertimage{liste_recus.png}{Tableau de tous les recus emis par le
president (dans l'onglet \textbf{Recus}) : tri par date,
beneficiaire, montant. Possibilite de telecharger chaque recu en
PDF ou de l'annuler.}

\insertimage{appel_president.png}{Onglet \textbf{Video} de
\code{PresidentDashboardComponent} : interface de programmation
d'appel video (LiveKit). Selection des participants (cible =
joueurs et entraineurs de SA discipline), date, duree, description.
Liste des appels a venir et passes avec jeton JWT signe 2h.}

\subsection{Back-office administrateur (29 sous-pages)}

Le dashboard admin (\code{DashboardComponent} route \code{/admin})
regroupe en sidebar 29 sous-pages \code{/admin/*} documentees
ci-dessous.

\insertimage{admin_demandes.png}{\code{AdminDemandesComponent}
(route \code{/admin/demandes}). Gestion des demandes de comptes
EN\_ATTENTE (Parent, Joueur, Entraineur, Staff, Journaliste,
President) avec motif, date, pieces KYC jointes, et boutons
Valider / Refuser.}

\insertimage{admin_presse_demandes.png}{\code{AdminPresseDemandesComponent}
(route \code{/admin/presse/demandes}). Onglet dedie aux demandes
d'accreditation presse (B.17) avec liste complete, motif
obligatoire au refus (@Valid size >= 5), et historique.}

\insertimage{admin_rapports.png}{\code{AdminRapportsComponent}
(route \code{/admin/rapports-financiers}). Rapports financiers
CSV / PDF avec filtres par periode, type (don / abonnement /
billet / boutique).}

\insertimage{admin_users.png}{\code{AdminUsersComponent} (route
\code{/admin/utilisateurs}). Liste complete des utilisateurs avec
filtres (role, statut, recherche). Actions : modifier le role,
activer / desactiver un compte, voir les details, supprimer (avec
garde-fou last-admin).}

\insertimage{admin_actualites.png}{\code{AdminActualitesComponent}
(route \code{/admin/actualites}). CRUD complet sur les actualites
du club (titre, contenu enrichi, image hero, categorie).}

\insertimage{admin_matchs.png}{\code{AdminMatchsComponent} (route
\code{/admin/matchs}). CRUD des matchs du calendrier officiel
(adversaire, date, competition, score, statut).}

\insertimage{admin_convocations.png}{\code{AdminConvocationsComponent}
(route \code{/admin/convocations}). Vue admin des convocations
(match + seance) creees par les entraineurs, avec possibilite de
valider ou refuser.}

\insertimage{admin_classements.png}{\code{AdminClassementsComponent}
(route \code{/admin/classements}). Mise a jour manuelle du
classement Botola Pro (points, victoires, nuls, defaites).}

\insertimage{admin_boutique.png}{\code{AdminBoutiqueComponent} (route
\code{/admin/boutique}). CRUD produits, categories, variantes,
images, et personnalisation maillot (numero + nom).}

\insertimage{admin_commandes.png}{\code{AdminCommandesComponent}
(route \code{/admin/commandes}). Suivi des commandes boutique :
statut (EN\_ATTENTE / PAYEE / EXPEDIEE / LIVREE), expedition,
facture PDF.}

\insertimage{admin_achats.png}{\code{AdminAchatsComponent} (route
\code{/admin/achats}). Onglets filtres pour les achats (dons,
abonnements, billets, boutique) avec export CSV / PDF.}

\insertimage{admin_billetterie.png}{\code{AdminBilletterieComponent}
(route \code{/admin/billetterie}). CRUD evenements + sections
(grille tarifaire TRIBUNE / VIP / ULTRA / VIRAGE) avec
PATCH /sections/\{id\} in-place pour modifier le prix.}

\insertimage{admin_plans_abonnement.png}{\code{AdminSubscriptionPlansComponent}
(route \code{/admin/abonnements/plans}). CRUD complet des plans
d'abonnement saisonnier avec upload d'une photo de carte
(Cloudinary) et champ API tier (standard / gold).}

\insertimage{admin_effectif.png}{\code{AdminEffectifComponent} (route
\code{/admin/effectif}). CRUD des fiches roster (joueurs + staff)
par discipline et categorie.}

\insertimage{admin_effectif_public.png}{\code{AdminJoueursPublicComponent}
(route \code{/admin/effectif-public}). Vue publique de l'effectif
(joueurs + staff), avec photo, position, categorie. Alimentee
automatiquement par l'effectif admin.}

\insertimage{admin_palmares.png}{\code{AdminPalmaresComponent} (route
\code{/admin/palmares}). CRUD des palmares avec upload photo
(via \code{api.uploadMedia()} + \code{api.getMediaUrl()}) ou URL
externe.}

\insertimage{admin_legendes.png}{\code{AdminLegendesComponent} (route
\code{/admin/legendes}). CRUD des legendes du club avec photo,
periode, palmares individuel.}

\insertimage{admin_seances.png}{\code{AdminSeancesComponent} (route
\code{/admin/seances}). Vue admin des seances d'entrainement
(lecture seule, planning global).}

\insertimage{admin_staff.png}{\code{AdminStaffComponent} (route
\code{/admin/staff}). CRUD des fiches staff (manager, preparateur
physique, medecin, etc.) par categorie.}

\insertimage{admin_academie.png}{\code{AdminAcademieComponent} (route
\code{/admin/academie} ou \code{/staff/academie}). CRUD des enfants
academie (inscription, validation, retrait).}

\insertimage{admin_mediatheque.png}{\code{AdminMediathequeComponent}
(route \code{/admin/mediatheque}). CRUD des medias partages dans
la mediateque tactique (videos, schemas, notes) avec visibilite
globale / par categorie / par joueur.}

\insertimage{admin_notifications.png}{\code{AdminNotificationsComponent}
(route \code{/admin/notifications}). Historique des notifications
in-app envoyees et planification de notifications groupees.}

\insertimage{admin_reclamations.png}{\code{AdminReclamationsComponent}
(route \code{/admin/reclamations}). Suivi des reclamations
utilisateurs avec reponse admin et cloture.}

\insertimage{admin_sondages.png}{\code{AdminSondagesComponent} (route
\code{/admin/sondages}). CRUD des sondages (choix multiples,
date de fermeture, publication).}

\insertimage{admin_elections.png}{\code{AdminElectionsComponent}
(route \code{/admin/elections}). CRUD des elections B.8 : titre,
candidats, dates, ouverture / fermeture / publication /
depublication (B.8.b), avec garde-fous fail-closed
(NOT\_ALL\_VOTED, PUBLISH\_NO\_ELIGIBLE\_SNAPSHOT).}

\insertimage{admin_badges.png}{\code{AdminBadgesComponent} (route
\code{/admin/badges}). CRUD des badges gamification (code
WAC\_*, description, image, condition d'obtention).}

\insertimage{admin_defis.png}{\code{AdminDefisComponent} (route
\code{/admin/defis}). CRUD des defis gamification C2 (type
SUBMIT\_PREDICTIONS / WIN\_EXACT\_SCORES / REACH\_POINTS, recompense,
periode).}

\insertimage{admin_bareme_points.png}{\code{AdminPointsRulesComponent}
(route \code{/admin/bareme-points}). Configuration du bareme de
points C3 (bonus prediction, score exact, signe, niveau) avec
rechargement du cache memoire.}

\insertimage{admin_parametres.png}{\code{AdminSettingsComponent}
(route \code{/admin/parametres}). Configuration globale de la
plateforme : sponsors (upload + redim thumbnailator), reseaux
sociaux (SVG + URL), backgrounds multi-image carrousel par page
(HOME / BILLETTERIE / ACADEMIE / ...), stade cover.}

\insertimage{not_found.png}{\code{NotFoundComponent} (route
\code{/**}). Page 404 avec recherche rapide et retour a
l'accueil.}

\section{Conclusion}

Ce chapitre a presente la phase de realisation structuree par couche
technique : backend Spring Boot 3.3, frontend Angular 19,
conteneurisation Docker, pipeline CI/CD GitHub Actions, et
deploiement Azure. La plateforme est aujourd'hui operationnelle,
securisee, modulaire et deployee en production.

% =====================================================================
% CONCLUSION GENERALE
% =====================================================================
\chapterbanner{CONCLUSION GENERALE}
\chapter*{Conclusion générale}
\addcontentsline{toc}{chapter}{Conclusion générale}

\section*{Bilan du projet}

Le projet de fin d'annee \textbf{Wydad Digital} avait pour objectif
la conception, le developpement et le deploiement d'une plateforme
numerique complete au service du Wydad Athletic Club. Realise sur
la periode \textbf{1\textsuperscript{er} -- 31 aout 2026}, le projet
a livre une solution operationnelle, securisee et modulaire couvrant
l'integralite des besoins exprimes par le porteur de projet.

\section*{Atteinte des objectifs}

Les objectifs initiaux ont ete atteints. La plateforme est fonctionnelle
avec 11 microservices couvrant les six piliers metier (adhesion et
billetterie, academie et roles sportifs, gouvernance participative,
communication officielle et interne, engagement communautaire,
administration) plus le cas C.21 (president de club). La securite
repose sur l'authentification JWT + refresh rotatif, la validation
cote serveur systematique, la gestion fine des 8 roles
utilisateurs, le secret partage inter-services a temps constant,
l'anti-IDOR strict, et l'User shadow pour les enfants. L'architecture
microservices est deployee de maniere reproductible sur une VM unique.
La VM Azure Standard\_B2s\_v2 (2 vCPU, 8~Go RAM, 30~Go SSD) tient dans
le credit Azure for Students.

\section*{Apports personnels}

Ce PFA m'a permis de mettre en pratique les connaissances acquises a
l'EMSI en architecture logicielle, genie logiciel, bases de donnees et
genie DevOps. J'ai decouvert et approfondi des technologies
professionnelles : Spring Boot 3, Angular 19, Docker, GitHub
Actions, JWT, WebSocket STOMP, LiveKit, Cloudinary, PostgreSQL.
J'ai conduit un projet de bout en bout en autonomie dans la periode
impartie. J'ai acquis une experience concrete du deploiement
d'applications sur le cloud Azure.

\section*{Limites et difficultes rencontrees}

La \textbf{contrainte de ressources VM} a ete un frein : 1~Go puis
8~Go RAM. La \textbf{securisation} a necessite plusieurs iterations
pour aboutir a un systeme durci. Le \textbf{cas B.18} (parent
supporter) a demande trois iterations de specification et de
refonte de l'experience utilisateur. Le \textbf{cas C.21} (president
de club) a necessite un backend nouveau et une refonte du check
messagerie.

\section*{Perspectives d'evolution}

Plusieurs pistes d'amelioration : internationalisation complete
(arabe, anglais) avec detection automatique de la langue ; application
mobile native (Ionic / Flutter / PWA installable) ; migration vers
Kubernetes (AKS) ; module e-cash / paiement integre avec un PSP
marocain (CMI, PayZone) ; dashboard analytique ; HTTPS systematique
via reverse proxy Caddy + Let's Encrypt ; recherche full-text ;
partie live du match (streaming) ; statistiques avancees pour le
staff sportif.

\section*{Mot de la fin}

Ce projet fut une experience formatrice complete : il m'a confronte
aux realites du metier d'ingenieur logiciel (qualite, securite,
dette technique, dialogue avec le client) et m'a appris a livrer un
produit fini, deploye et documente, dans le respect des bonnes
pratiques. Je suis convaincu que la plateforme \textit{Wydad Digital}
constitue une base solide et evolutive au service de la communaute
du Wydad Athletic Club.

% =====================================================================
% BIBLIOGRAPHIE
% =====================================================================
\chapterbanner{BIBLIOGRAPHIE ET NETOGRAPHIE}
\chapter*{Bibliographie et Nétographie}
\addcontentsline{toc}{chapter}{Bibliographie et Nétographie}

\section*{Ouvrages et supports de cours}

\begin{itemize}
  \item \textbf{Spring Documentation} -- \textit{Reference
        Documentation}, VMware / Broadcom, 2024.
        \url{https://docs.spring.io/spring-boot/docs/current/reference/html/}
  \item \textbf{Spring Security} -- \textit{Reference Documentation},
        VMware / Broadcom, 2024.
        \url{https://docs.spring.io/spring-security/reference/}
  \item \textbf{Angular Documentation} -- Google / Angular Team, 2024.
        \url{https://angular.dev/}
  \item Fowler M., \textit{Patterns of Enterprise Application
        Architecture}, Addison-Wesley, 2002.
  \item Newman S., \textit{Building Microservices}, O'Reilly Media,
        2\textsuperscript{ème} édition, 2021.
  \item Richardson C., \textit{Microservices Patterns}, Manning
        Publications, 2018.
  \item Hunt A., Thomas D., \textit{The Pragmatic Programmer},
        Addison-Wesley, 2019.
  \item Bass L., Weber I., Zhu L., \textit{DevOps: A Software
        Architect's Perspective}, Addison-Wesley, 2015.
  \item Humble J., Farley D., \textit{Continuous Delivery},
        Addison-Wesley, 2010.
\end{itemize}

\section*{Articles et standards}

\begin{itemize}
  \item RFC 7519 -- \textit{JSON Web Token (JWT)}.
        \url{https://datatracker.ietf.org/doc/html/rfc7519}
  \item RFC 6455 -- \textit{The WebSocket Protocol}.
        \url{https://datatracker.ietf.org/doc/html/rfc6455}
  \item OWASP Top 10 -- \textit{Web Application Security Risks},
        2021.
        \url{https://owasp.org/Top10/}
  \item Docker Documentation -- \textit{Get Started, Compose
        Specification}.
        \url{https://docs.docker.com/}
  \item Microsoft Azure -- \textit{Virtual Machines documentation}.
        \url{https://learn.microsoft.com/azure/virtual-machines/}
  \item LiveKit Documentation -- \textit{Self-hosted WebRTC}.
        \url{https://docs.livekit.io/}
  \item Cloudinary Documentation -- \textit{Media management
        platform}.
        \url{https://cloudinary.com/documentation}
\end{itemize}

\section*{Netographie (sites consultes)}

\begin{itemize}
  \item Site officiel du Wydad Athletic Club :
        \url{https://www.wac.ma/}
  \item Spring Initializr :
        \url{https://start.spring.io/}
  \item Stack Overflow :
        \url{https://stackoverflow.com/}
  \item Baeldung :
        \url{https://www.baeldung.com/}
\end{itemize}

% =====================================================================
% ANNEXES
% =====================================================================
\chapterbanner{ANNEXES}
\chapter*{Annexes}
\addcontentsline{toc}{chapter}{Annexes}

\section*{Annexe A : Liste des diagrammes}

Ce rapport integre les diagrammes suivants, generes a partir du
code source du projet (fichiers \code{.puml} dans
\code{docs/diagrams/final/}) via PlantUML 1.2024.7, conformement
au code qui compile (controllers, services, DTOs back Spring Boot
3.3 + composants standalone Angular 19 front) :

\begin{itemize}
  \item \textbf{12 diagrammes de cas d'utilisation} (UML 2.5) :
    \begin{itemize}
      \item UC-01a vue globale supporter (VISITEUR, ADHERENT,
            PARENT, JOURNALISTE)
      \item UC-01b vue globale club (JOUEUR, ENTRAINEUR, STAFF,
            PRESIDENT, ADMIN)
      \item UC-02 supporter public (consultation libre)
      \item UC-03 supporter achat (billetterie + abonnement +
            boutique)
      \item UC-04 supporter academie (B.18 parent supporter)
      \item UC-05 supporter gamification (C1/C2/C3)
      \item UC-06 supporter votes (elections B.8 + sondages)
      \item \textbf{UC-07 joueur} (VAGUE 2 enrichi, 25 UCs, 8
            packages : B.12, B.18.b, B.28, V2.5, Phase 0)
      \item \textbf{UC-08 staff + entraineur} (VAGUE 2 enrichi,
            heritage, 24 UCs, 8 packages)
      \item \textbf{UC-09 president} (VAGUE 2 enrichi, 22 UCs,
            8 packages : C.21 recus, V2.5, B.12)
      \item \textbf{UC-10 messagerie unifiee V2.5} (4 acteurs,
            16 UCs transversaux, team-chat + 1-to-1 + canal
            president + selecteur libre)
      \item \textbf{UC-11 admin} (VAGUE 2 enrichi, 42 UCs,
            7 packages, 29 sous-pages /admin/*)
      \item UC-12 journaliste (B.17 accreditation)
    \end{itemize}
  \item \textbf{10 diagrammes de sequence} :
    sequence-achat-billet-cas1-visiteur (B.28),
    sequence-achat-billet-cas1-paiement,
    sequence-achat-billet-cas2-parent (B.18),
    sequence-achat-billet-cas2-parent-billet (B.18),
    sequence-vote-election (B.8),
    sequence-vote-election-failures (B.8),
    sequence-accreditation-presse (B.17),
    sequence-convocations (B.18.b),
    sequence-feuille-match (B.18.b),
    sequence-team-chat (V2.2 STOMP),
    sequence-president-receipts (C.21),
    sequence-parent-academie (B.18),
    sequence-admin-validation-comptes (Phase 0),
    sequence-abonnement (B.12).
  \item \textbf{1 diagramme de collaboration} : collab-accreditation
    (B.17, focus sur l'ordre temporel des messages).
  \item \textbf{3 diagrammes d'etats-transition} (cycles de vie
    metier du code) :
    \begin{itemize}
      \item state-ticket (cycle de vie Billet, 5 etats)
      \item \textbf{state-election} (VAGUE 2 enrichi, B.8 / B.8.b
            / B.8.c / B.8.e / B.8.h, 3 etats, scheduler 60s,
            fail-closed)
      \item \textbf{state-press-accred} (VAGUE 2 enrichi, B.17,
            3 etats, IDOR badge strict, motif obligatoire,
            pas de re-soumission)
    \end{itemize}
  \item \textbf{8 diagrammes de classes} (vue macro + 7 par
    domaine) : classes-macro, classes-auth, classes-billetterie,
    classes-sports, classes-shop, classes-academie (B.18),
    classes-messagerie (V2.5), classes-election (B.8),
    classes-gamification (C1/C2/C3).
  \item \textbf{4 schemas d'architecture} (non-PlantUML, schemas
    structurels) : architecture globale microservices,
    architecture Docker, pipeline CI/CD GitHub Actions,
    deploiement Azure VM B2s\_v2.
\end{itemize}

\paragraph{Principes de modelisation respectes.}
La modelisation suit strictement les conventions UML 2.5 :
\begin{itemize}
  \item Pas de chevauchement entre packages (chaque UC est rattache
        a un seul sous-domaine fonctionnel).
  \item Un UC par action metier (atomicite des cas d'utilisation).
  \item Swimlanes par acteur (un diagramme par profil, ou separation
        claire en swimlanes sur les diagrammes transversaux comme
        UC-10 messagerie unifiee).
  \item Relations \code{<<include>>} losange plein uniquement pour
        les dependances obligatoires (ex. login).
  \item Relations \code{<<extend>>} losange creux uniquement pour
        les comportements optionnels (ex. upload photo).
  \item Diagrammes issus du code qui compile : pas d'invention, pas
        d'idealisation, uniquement ce qui est implementable
        aujourd'hui.
\end{itemize}

\section*{Annexe B : Glossaire}

\begin{itemize}
  \item \textbf{BCrypt} : algorithme de hashage de mots de passe avec
        un cout configurable.
  \item \textbf{Cloudinary} : service tiers de gestion de medias
        (images, videos) avec transformation a la volee.
  \item \textbf{Composed mapping} : mecanisme Spring Data qui permet
        de definir plusieurs mappings JPA pour la meme entite.
  \item \textbf{Docker Compose} : outil de definition et d'execution
        d'applications Docker multi-conteneurs.
  \item \textbf{HMAC-SHA256} : algorithme de signature numerique
        utilisant une cle secrete et une fonction de hashage SHA-256.
  \item \textbf{JWT (JSON Web Token)} : standard ouvert (RFC 7519) pour
        la transmission securisee d'informations entre parties.
  \item \textbf{LiveKit} : plateforme WebRTC open-source pour les
        appels video et audio en temps reel.
  \item \textbf{Monorepo} : strategie de gestion de code source ou
        plusieurs projets partagent un meme depot Git.
  \item \textbf{STOMP (Simple Text Oriented Messaging Protocol)} :
        protocole de messagerie simple au-dessus de WebSocket.
  \item \textbf{WebSocket} : protocole de communication full-duplex
        sur une seule connexion TCP (RFC 6455).
\end{itemize}

\end{document}
